% ACM Algorithm Template - v4.0
% Yume Maruyama <kirainmoe@gmail.com>

\documentclass[12pt]{report}

\usepackage{fontspec}
\usepackage{graphicx}
\usepackage{geometry}
\usepackage{listings}
\usepackage[dvipsnames, svgnames, x11names]{xcolor}
\usepackage{fontspec}
\usepackage{mathtools}
\usepackage{multicol}
\usepackage{verbatim} 
\usepackage{amsthm}
\usepackage{amssymb}
\usepackage{fancyhdr}
\usepackage{minted}
\usepackage{verbatim}
\usepackage{pdfpages}
\usepackage{setspace}  %设置行间距
\usepackage[colorlinks=true]{hyperref}
\usepackage[BoldFont,SlantFont,CJKchecksingle]{xeCJK}
\usepackage{tcolorbox}
\usepackage{lastpage} %总页数
\usepackage{fancyhdr} %使用fancyhdr
\usepackage{indentfirst}

% 设置字体
\setCJKmainfont{SimSun}
\setmainfont{Times New Roman}
\setsansfont{Fira Code}
% 页边距，行距设置
\geometry{left=1.5cm,right=1.5cm,top=1.5cm,bottom=1.5cm}
\linespread{1.2}

% link 颜色
\hypersetup{
    colorlinks=true,
    linkcolor=black,
}

% listings 样式
\lstset{
    basicstyle = \ttfamily,
}

% 代码样式，使用 minted
\setminted[c++]{linenos, frame=lines, framesep=2mm, baselinestretch=1.0, breaklines, mathescape, numbers=left, numbersep=2mm}
\setminted[java]{linenos, frame=leftline, framesep=2mm, baselinestretch=1.0, breaklines, mathescape, bgcolor=LavenderBlush1}
\setminted[python]{linenos, frame=leftline, framesep=2mm, baselinestretch=1.0, breaklines, mathescape, bgcolor=LavenderBlush1}
\setminted[batch]{linenos, frame=leftline, framesep=2mm, baselinestretch=1.0, breaklines, mathescape, bgcolor=LavenderBlush1}
\setminted[bash]{linenos, frame=leftline, framesep=2mm, baselinestretch=1.0, breaklines, mathescape, bgcolor=LavenderBlush1}
\setminted[json]{linenos, frame=leftline, framesep=2mm, baselinestretch=1.0, breaklines, mathescape, bgcolor=LavenderBlush1}
\usemintedstyle{vs}
\definecolor{primary-blue}{RGB}{9,109,217}
\definecolor{cyan}{RGB}{19,194,194}

% 标题 & 作者信息
\title{\Huge \textbf{ICPC Templates} \linebreak \linebreak \large \textcolor{primary-blue}{\textbf{我们需要更深入浅出一些}}}
\author{Mine\_King, Tx\_Lcy, 369Pai \\ \small 洛谷科技}

% 设置页眉风格
\pagestyle{fancy}                         %设置页眉页脚
\lhead{}                                  %页眉左侧显示页数                  
\chead{}                                  %页眉中
\rhead{\small\leftmark}                   %章节信息                       
\cfoot{\thepage}                          %当前页，记得调用前文提到的宏包 
\rfoot{}%                                 %页脚左               
\lfoot{}%					          %页脚右
\renewcommand{\headrulewidth}{0mm}        %页眉线宽，设为0可以去页眉线
\renewcommand{\footrulewidth}{0mm}        %页脚线宽，设为0可以去页脚线
\renewcommand*\contentsname{目录}

\begin{document}
% 用\setlength设置页眉页脚边距
%\setlength{\voffset}{-10mm}                        
%\setlength{\topmargin}{0mm}
\setlength{\headheight}{5mm}
\setlength{\headsep}{5mm}
%\setlength{\footskip}{10mm}
%\setlength{\baselineskip}{20pt}  %设置行间距

% 标题页面
\begin{titlepage}
    \begin{figure}
        \centering
        \includegraphics[height=8cm]{images/logo.jpg}
        % \par\includegraphics[height=3cm]{images/acm.jpg}
    \end{figure}
    
    \maketitle
\end{titlepage}

%\clearpage
%\phantom{s}
%\setcounter{page}{0}
%\thispagestyle{empty}
%\clearpage

% 目录
\pagenumbering{gobble}
\tableofcontents       % 生成目录
\clearpage

\pagenumbering{arabic}

\newpage
\setcounter{page}{1}   % 从下面开始编页码

% 正文内容

%% Chapter ? 字符串
\chapter{字符串}

\section{最小表示法}

\begin{minted}{c++}
int i = 0, j = 1, k = 0;
while (k < n && i < n && j < n)
    if (a[(i + k) % n] == a[(j + k) % n]) k++;
    else {
        if (a[(i + k) % n] > a[(j + k) % n]) i = i + k + 1;
        else j = j + k + 1;
        if (i == j) i++;
        k = 0;
    }
ans = min(i, j);
\end{minted}

\section{Border 理论}

\subsection{关键结论}

\begin{tcolorbox}
\textbf{定理 1}：对于一个字符串 $s$，若用 $t$ 表示其最长的 Border，则有 $\mathcal{B}(s) = \mathcal{B}(t) \cup \{t\}$。
\end{tcolorbox}

\begin{tcolorbox}
\textbf{定理 2}：一个字符串的 Border 与 Period 一一对应。具体地，$\mathrm{pre}(s, i) \in \mathcal{B}(s) \iff |s| - i \in \mathcal{P}(s)$。
\end{tcolorbox}

\begin{tcolorbox}
\textbf{弱周期引理}：

$$
\forall p, q\in\mathcal{P}(s), p + q \le |s| \implies \gcd(p, q)\in\mathcal{P}(s)
$$
\end{tcolorbox}

\begin{tcolorbox}
\textbf{定理 3}：若字符串 $t$ 是字符串 $s$ 的前缀，且 $a \in \mathcal{P}(s), b \in \mathcal{P}(t), b \mid a, |t| \ge a$，且 $b$ 是 $t$ 的整周期，则有 $b \in \mathcal{P}(s)$。
\end{tcolorbox}

\begin{tcolorbox}
\textbf{周期引理}：

$$
\forall p, q\in\mathcal{P}(s), p + q - \gcd(p, q) \le |s| \implies \gcd(p, q)\in\mathcal{P}(s)
$$
\end{tcolorbox}

\begin{tcolorbox}
\textbf{定理 4}：对于文本串 $s$ 和模式串 $t$，若 $|t| \ge \frac{|s|}{2}$，且 $t$ 在 $s$ 中至少成功匹配了 $3$ 次，则每次匹配的位置形成一个等差数列，且公差为 $t$ 的最小周期。
\end{tcolorbox}

\begin{tcolorbox}
\textbf{定理 5}：一个字符串 $s$ 的所有长度不小于 $\frac{|s|}{2}$ 的 Border 的长度构成一个等差数列。
\end{tcolorbox}

\begin{tcolorbox}
\textbf{定理 6}：一个字符串的所有 Border 的长度排序后可以划分成 $\lceil\log_2|s|\rceil$ 个连续段，使得每段都是一个等差数列。
\end{tcolorbox}

\begin{tcolorbox}
\textbf{定理 7}：回文串的回文前/后缀即为该串的 Border。
\end{tcolorbox}

\begin{tcolorbox}
\textbf{定理 8}：若回文串 $s$ 有周期 $p$，则可以把 $\mathrm{pre}(s, p)$ 划分成长度为 $|s| \bmod p$ 的前缀和长度为 $p - |s| \bmod p$ 的后缀，使得它们都是回文串。
\end{tcolorbox}

\begin{tcolorbox}
\textbf{定理 9}：若 $t$ 是回文串 $s$ 的最长 Border 且 $|t| \ge \frac{|s|}{2}$，则 $t$ 在 $s$ 中只能匹配 $2$ 次。
\end{tcolorbox}

\begin{tcolorbox}
\textbf{定理 10}：对于任意一个字符串以及 $u, v\in\mathrm{Ssuf}(s), |u| < |v|$，一定有 $u$ 是 $v$ 的 Border。
\end{tcolorbox}

\begin{tcolorbox}
\textbf{定理 11}：对于任意一个字符串 $s$ 以及 $u, v\in\mathrm{Ssuf}(s), |u| < |v|$，一定有 $2|u| \le |v|$。
\end{tcolorbox}

\begin{tcolorbox}
\textbf{定理 12}：$|\mathrm{Ssuf}(s)| \le \log_2|s|$。
\end{tcolorbox}

\subsection{KMP}

\begin{minted}{c++}
// n is |s|, m is |t|.
for (int i = 1, j = 0; i <= n; i++) {
    while (j && t[j + 1] != s[i]) j = pi[j];
    if (t[j + 1] == s[i])
        if (++j == m) {
            // ...
            j = nxt[j];
        }
}
\end{minted}

\section{Z 函数}

\begin{minted}{c++}
// n is |s|.
for (int i = 2, j = 0; i <= n; i++) {
    if (i < j + z[j]) z[i] = min(z[i - j + 1], j + z[j] - i);
    while (i + z[i] <= n && s[i + z[i]] == s[1 + z[i]]) z[i]++;
    if (i + z[i] > j + z[j]) j = i;
}
\end{minted}

\section{Manacher}

\begin{minted}{c++}
// == Main ==
// t is the original string, n is |t|.
string s = "^#";
for (char i : t) s.push_back(i), s.push_back('#');
s.push_back ('@');
for (int i = 1, j = 0; i <= 2 * n + 1; i++) {
    if (i <= j + p[j]) p[i] = min (p[2 * j - i], j + p[j] - i);
    while (s[i - p[i] - 1] == s[i + p[i] + 1]) p[i]++;
    if (i + p[i] > j + p[j]) j = i;
}
\end{minted}

\section{AC 自动机}

\begin{minted}{c++}
// == Preparations ==
#include <queue>
// == Main ==
struct ACAM {
    int tot, fail[200005], delta[200005][26];

    void insert(string s, int id) {
        int now = 0;
        for (char c : s) {
            int v = c - 'a';
            if (!delta[now][v]) delta[now][v] = ++tot;
            now = delta[now][v];
        }
        return;
    }
    void build() {
        queue<int> q;
        for (int c = 0; c < 26; c++)
            if (delta[0][c]) q.push(delta[0][c]);
        while (!q.empty()) {
            int now = q.front();
            q.pop();
            for (int c = 0; c < 26; c++)
                if (delta[now][c]) fail[delta[now][c]] = delta[fail[now]][c], q.push(delta[now][c]);
                else delta[now][c] = delta[fail[now]][c];
        }
        return;
    }
} ac;
\end{minted}

\section{回文自动机 PAM}

\begin{minted}{c++}
// == Main ==
struct PAM {
    int tot, delta[500005][26], len[500005], fail[500005], ans[500005];
    string s;
    int lst;

    PAM() {tot = 1; len[0] = 0; len[1] = -1; fail[0] = fail[1] = 1;}
    int getfail(int now, int i) {
        while (s[i - len[now] - 1] != s[i]) now = fail[now];
        return now;
    }
    void insert(int i) {
        int now = getfail(lst, i);
        if (!delta[now][s[i] - 'a']) {
            len[++tot] = len[now] + 2;
            fail[tot] = delta[getfail(fail[now], i)][s[i] - 'a'];
            delta[now][s[i] - 'a'] = tot;
            ans[tot] = ans[fail[tot]] + 1;
        }
        lst = delta[now][s[i] - 'a'];
        return;
    }
} p;
\end{minted}

\section{后缀自动机}

\subsection{普通 SAM}

\begin{minted}{c++}
// == Main ==
struct SAM {
    int tot, lst;
    int len[2000005], siz[2000005], link[2000005];
    int delta[2000005][26];

    SAM() {link[0] = -1;}
    void insert(char ch) {
        int c = ch - 'a', now = ++tot;
        len[now] = len[lst] + 1;
        siz[now] = 1;
        for (int p = lst; p != -1; p = link[p])
            if (!delta[p][c]) delta[p][c] = tot;
            else if (len[delta[p][c]] == len[p] + 1) {link[now] = delta[p][c]; break;}
            else {
                int q = delta[p][c], v = ++tot;
                len[v] = len[p] + 1;
                memcpy(delta[v], delta[q], sizeof(delta[v]));
                link[v] = link[q], link[q] = v, link[now] = v;
                for (int i = p; delta[i][c] == q; i = link[i]) delta[i][c] = v;
                break;
            }
        lst = now;
        return ;
    }
} sam;
\end{minted}

\subsection{广义 SAM}

注意自动机空间要开 Trie 的两倍。

\begin{minted}{c++}
struct GSAM {
    int tot;
    int delta[2000005][26], link[2000005], len[2000005];
    struct Trie {
        int tot, trie[1000005][26], st[1000005];

        void insert(string s) {
            int now = 0;
            for (char c : s) {
                int id = c - 'a';
                if (!trie[now][id]) trie[now][id] = ++tot;
                now = trie[now][id];
            }
            return;
        }
    } tr;

    GSAM() {link[0] = -1;}
    int insert(int c, int lst) {
        int now = ++tot;
        len[now] = len[lst] + 1;
        for (int p = lst; p != -1; p = link[p])
            if (!delta[p][c]) delta[p][c] = now;
            else if (len[delta[p][c]] == len[p] + 1) {link[now] = delta[p][c]; break;}
            else {
                int q = delta[p][c], v = ++tot;
                len[v] = len[p] + 1;
                memcpy(delta[v], delta[q], sizeof(delta[v]));
                link[v] = link[q], link[q] = v, link[now] = v;
                for (int i = p; i != -1 && delta[i][c] == q; i = link[i]) delta[i][c] = v;
                break;
            }
        return now;
    }
    void build() {
        queue<int> q;
        tr.st[0] = 0;
        q.push(0);
        while (!q.empty()) {
            int now = q.front();
            q.pop();
            for (int i = 0; i < 26; i++)
                if (tr.trie[now][i])
                    tr.st[tr.trie[now][i]] = insert(i, tr.st[now]), q.push(tr.trie[now][i]);
        }
        return;
    }
} gsam;
\end{minted}

\section{后缀排序}

\begin{minted}{c++}
// == Preparations ==
int sa[2000005], rk[2000005], b[1000005], cp[2000005];
// == Main ==
for (int i = 1; i <= n; i++) b[rk[i] = s[i]]++;
for (int i = 1; i < 128; i++) b[i] += b[i - 1];
for (int i = n; i >= 1; i--) sa[b[rk[i]]--] = i;
memcpy(cp, rk, sizeof(cp));
for (int i = 1, j = 0; i <= n; i++)
    if (cp[sa[i]] == cp[sa[i - 1]]) rk[sa[i]] = j;
    else rk[sa[i]] = ++j;
for (int w = 1; w < n; w <<= 1) {
    memcpy(cp, sa, sizeof(cp));
    memset(b, 0, sizeof(b));
    for (int i = 1; i <= n; i++) b[rk[cp[i] + w]]++;
    for (int i = 1; i <= n; i++) b[i] += b[i - 1];
    for (int i = n; i >= 1; i--) sa[b[rk[cp[i] + w]]--] = cp[i];
    memcpy(cp, sa, sizeof(cp));
    memset(b, 0, sizeof(b));
    for (int i = 1; i <= n; i++) b[rk[cp[i]]]++;
    for (int i = 1; i <= n; i++) b[i] += b[i - 1];
    for (int i = n; i >= 1; i--) sa[b[rk[cp[i]]]--] = cp[i];
    memcpy(cp, rk, sizeof(cp));
    for (int i = 1, j = 0; i <= n; i++)
        if (cp[sa[i]] == cp[sa[i - 1]] && cp[sa[i] + w] == cp[sa[i - 1] + w]) rk[sa[i]] = j;
        else rk[sa[i]] = ++j;
}
\end{minted}

求 height：

\begin{minted}{c++}
for (int i = 1, k = 0; i <= n; i++) {
    if (k) k--;
    while (i + k <= n && sa[rk[i] - 1] + k <= n && s[i + k] == s[sa[rk[i] - 1] + k]) k++;
    height[rk[i]] = k;
}
\end{minted}

$$
LCP(sa_i, sa_j) = \min\limits_{k = i + 1}^j \{height_k\}
$$
%
%% Chapter ? 数据结构
\chapter{数据结构}

\section{K-D Tree}

\begin{minted}{c++}
// == Main ==
template<const int Dim = 2>
struct KDTree {
    using point = array<int, Dim>;
    struct node {
        point p, l, r;
        int val, siz, sum;
        node *ls, *rs;

        node() = default;
        node(point _p, int _val = 0):
            p(_p), l(_p), r(_p), val(_val), siz(1), sum(_val), ls(nullptr), rs(nullptr) {}
        void pushup() {
            l = r = p, siz = 1, sum = val;
            for (int i = 0; i < Dim; i++) {
                if (ls) l[i] = min(l[i], ls->l[i]), r[i] = max(r[i], ls->r[i]);
                if (rs) l[i] = min(l[i], rs->l[i]), r[i] = max(r[i], rs->r[i]);
            }
            if (ls) siz += ls->siz, sum += ls->sum;
            if (rs) siz += rs->siz, sum += rs->sum;
            return ;
        }
    };
    vector<node *> root;
    using itor = typename vector<node>::iterator;

    node *build(itor l, itor r, int dim = 0) {
        if (l == r) return nullptr;
        int mid = (r - l) / 2;
        nth_element(l, l + mid, r, [&dim](const node &x, const node &y) {return x.p[dim] < y.p[dim];});
        node *now = new node(*(l + mid));
        now->ls = build(l, l + mid, (dim + 1) % Dim);
        now->rs = build(l + mid + 1, r, (dim + 1) % Dim);
        now->pushup();
        return now;
    }
    void getnode(node *now, vector<node> &vec) {
        if (!now) return ;
        vec.push_back(*now);
        getnode(now->ls, vec), getnode(now->rs, vec);
        delete now;
        return ;
    }
    void insert(point p, int val) {
        vector<node> tmp({node(p, val)});
        while (!root.empty() && root.back()->siz == (int)tmp.size())
            getnode(root.back(), tmp), root.pop_back();
        sort(tmp.begin(), tmp.end(), [](const node &x, const node &y) {return x.p < y.p;});
        vector<node> vec;
        for (node i : tmp)
            if (!vec.empty() && vec.back().p == i.p) vec.back().val += i.val;
            else vec.push_back(i);
        root.push_back(build(vec.begin(), vec.end()));
        return ;
    }
    int query(point ll, point rr, node *now) {
        if (!now) return 0;
        int flag = 1;
        for (int i = 0; i < Dim; i++)
            if (now->r[i] < ll[i] || now->l[i] > rr[i]) return 0;
            else flag &= ll[i] <= now->l[i] && now->r[i] <= rr[i];
        if (flag) return now->sum;
        flag = 1;
        for (int i = 0; i < Dim; i++) flag &= ll[i] <= now->p[i] && now->p[i] <= rr[i];
        return flag * now->val + query(ll, rr, now->ls) + query(ll, rr, now->rs);
    }
    int query(point ll, point rr) {
        int ans = 0;
        for (node *rt : root) ans += query(ll, rr, rt);
        return ans;
    }
    ~KDTree() {
        vector<node> tmp;
        for (node *rt : root) getnode(rt, tmp);
    }
};
\end{minted}

\section{Link Cut Tree}

代码维护的是点权异或和。

\begin{minted}{c++}
// == Main ==
struct LinkCutTree {
    int fa[100005], son[100005][2], siz[100005], swp[100005];
    int val[100005], Xor[100005];

    void pushup(int now) {
        siz[now] = siz[son[now][0]] + siz[son[now][1]] + 1;
        Xor[now] = Xor[son[now][0]] ^ val[now] ^ Xor[son[now][1]]; // 此处更新信息。
        return;
    }
    void pushdown(int now) {
        if(!swp[now]) return ;
        swap(son[now][0], son[now][1]);
        swp[son[now][0]] ^= 1, swp[son[now][1]] ^= 1;
        swp[now] = 0;
        // 此处将信息 pushdown
        return;
    }
    int isRoot(int now) {return now != son[fa[now]][0] && now != son[fa[now]][1];}
    int get(int now) {return now == son[fa[now]][1];}
    void rotate(int x) {
        int y = fa[x], z = fa[fa[x]], chk = get(x);
        if (!isRoot(y)) son[z][get(y)] = x;
        son[y][chk] = son[x][chk ^ 1], fa[son[x][chk ^ 1]] = y;
        son[x][chk ^ 1] = y, fa[y] = x;
        fa[x] = z;
        pushup(y), pushup(x);
        return;
    }
    void splay(int now) {
        vector<int> stk;
        stk.push_back(now);
        for (int i = now; !isRoot(i); i = fa[i]) stk.push_back(fa[i]);
        while (!stk.empty()) pushdown(stk.back()), stk.pop_back();
        for (int f; f = fa[now], !isRoot(now); rotate(now))
            if (!isRoot(f)) rotate(get(f) == get(now) ? f : now);
        return;
    }
    void access(int now) { // 打通到根的链
        for (int lst = 0; now; lst = now, now = fa[now]) splay(now), son[now][1] = lst, pushup(now);
        return;
    }
    void makeRoot(int now) {access(now); splay(now); swp[now] ^= 1; return;} // 设置根
    void link(int u, int v) {makeRoot(u); fa[u] = v; return;} // 连接
    void cut(int u, int v) {makeRoot(u); access(v); splay(v); son[v][0] = fa[u] = 0; return;} // 切割
    int find(int now) { // 找根
        access(now), splay(now);
        pushdown(now);
        while (son[now][0]) now = son[now][0], pushdown(now);
        splay(now);
        return now;
    }
    void split(int u, int v) {makeRoot(u); access(v); splay(u); return;} // 剖出 u ~ v 的链
    void update(int u, int _val) {split(u, u); val[u] = Xor[u] = _val; return ;} // 修改操作，split 后做就行了，此处为单点修改。
    int query(int u, int v) {split(u, v); return Xor[u];} // 查询操作，split 后做就行了。
    int isConnected(int u, int v) {return find(u) == find(v);} // 查询两个点是否连通。
};
\end{minted}
%
%% Chapter ? 数论
\chapter{数论}

\section{Miller Rabin 和 Pollard Rho}

\subsection{Miller Rabin}

\begin{minted}{c++}
// == Preparations ==
const int prime[] = {2, 3, 5, 7, 9, 11, 13, 17, 19, 23, 29, 31, 37};

long long power(long long a, long long b, long long mod) {
    long long ans = 1;
    while (b) {
        if (b & 1) ans = (__int128)ans * a % mod;
        a = (__int128)a * a % mod;
        b >>= 1;
    }
    return ans % mod;
}
// == Main ==
inline int Miller_Rabin(long long n) {
    if (n == 1) return 0;
    if (n == 2) return 1;
    if (n % 2 == 0) return 0;
    long long u = n - 1, t = 0;
    while (u % 2 == 0) u /= 2, t++;
    for (int i = 0; i < 12; i++) {
        if (prime[i] % n == 0) continue;
        long long x = power(prime[i] % n, u, n);
        if (x == 1) continue;
        int flag = 0;
        for (int j = 1; j <= t; j++) {
            if (x == n - 1) {flag = 1; break;}
            x = (__int128)x * x % n;
        }
        if (!flag) return 0;
    }
    return 1;
}
\end{minted}

\subsection{Pollard Rho}

\begin{minted}{c++}
// == Preparations ==
#include <chrono>
#include <random>

mt19937_64 gen(chrono::system_clock::now().time_since_epoch().count());
/*
Miller Rabin
*/
// == Main ==
long long Pollard_Rho(long long n) {
    long long s = 0, t = 0, c = gen() % (n - 1) + 1;
    for (int goal = 1; ; goal <<= 1, s = t) {
        long long val = 1;
        for (int step = 1; step <= goal; step++) {
            t = ((__int128)t * t + c) % n;
            val = (__int128)val * abs(t - s) % n;
            if (!val) return n;
            if (step % 127 == 0) {
                long long d = __gcd(val, n);
                if (d > 1) return d;
            }
        }
        long long d = __gcd(val, n);
        if (d > 1) return d;
    }
}
void factor(long long n) {
    if (n < ans) return ;
    if (n == 1 || Miller_Rabin(n)) {
        // n = 1 或 n 是一个质因子。
        // ...
    }
    long long p;
    do p = Pollard_Rho(n);
    while (p == n);
    while (n % p == 0) n /= p;
    factor(n), factor(p);
    return ;
}
\end{minted}

\section{常用数论算法}

\subsection{exgcd}

求出来的解满足 $|x| \le b, |y| \le a$。

\begin{minted}{c++}
// == Main ==
int exgcd(int a, int b, int &x, int &y) {
    if (b == 0) {x = 1; y = 0; return a;}
    int m = exgcd(b, a % b, y, x);
    y -= a / b * x;
    return m;
}
\end{minted}

\subsection{CRT}

\textbf{没用。}

\begin{minted}{c++}
// == Preparations ==
#include <vector>
// == Main ==
int CRT(vector<pair<int, int>> &a) {
    int M = 1;
    for (auto i : a) M *= i.second;
    int res = 0;
    for (auto i : a) {
        int t = inv(M / i.second, i.second);
        res = (res + (long long)i.first * (M / i.second) % M * t % M) % M;
    }
    return res;
}
\end{minted}

\subsection{exCRT}

需保证 $\mathrm{lcm}$ 在 \lstinline|long long| 范围内。

\begin{minted}{c++}
// == Preparations ==
long long exgcd(long long a, long long b, long long &x, long long &y);
// == Main ==
long long exCRT(vector<pair<long long, long long>> vec) {
    long long ans = vec[0].first, mod = vec[0].second;
    for (int i = 1; i < (int)vec.size(); i++) {
        long long a = mod, b = vec[i].second, c = vec[i].first - ans % b;
        long long x, y;
        long long g = exgcd(a, b, x, y);
        if (c % g != 0) return -1;
        b /= g;
        x = (__int128)x * (c / g) % b;
        ans += x * mod;
        mod *= b;
        ans = (ans % mod + mod) % mod;
    }
    return ans;
}
\end{minted}

\subsection{exLucas}

\begin{minted}{c++}
// == Preparations ==
int power(int a, long long b, int mod);
int exgcd(int a, int b, int &x, int &y);
int CRT(vector<pair<int, int>> &a);
// == Main ==
int inv(int n, int p) {
    int x, y;
    exgcd(n, p, x, y);
    return (x % p + p) % p;
}
int fac(long long n, int p, int pk) {
    if (n == 0) return 1;
    int res = 1;
    for (int i = 1; i < pk; i++)
        if (i % p != 0) res = (long long)res * i % pk;
    res = power(res, n / pk, pk);
    for (int i = 1; i <= n % pk; i++)
        if (i % p != 0) res = (long long)res * i % pk;
    return (long long)res * fac(n / p, p, pk) % pk;
}
int C(long long n, long long m, int p, int pk) {
    long long x = n, y = m, z = n - m;
    int res = (long long)fac(x, p, pk) * inv(fac(y, p, pk), pk) % pk * inv(fac(z, p, pk), pk) % pk;
    long long e = 0;
    while (x) e += x / p, x /= p;
    while (y) e -= y / p, y /= p;
    while (z) e -= z / p, z /= p;
    return (long long)res * power(p, e, pk) % pk;
}
int exLucas(long long n, long long m, int p) {
    vector<pair<int, int>> a;
    for (int i = 2; i * i <= p; i++)
        if (p % i == 0) {
            int pk = 1;
            while (p % i == 0) pk *= i, p /= i;
            a.emplace_back(C(n, m, i, pk), pk);
        }
    if (p != 1) a.emplace_back(C(n, m, p, p), p);
    return CRT(a);
}
\end{minted}

\section{万能欧几里得算法}

\begin{tcolorbox}
问题描述：\\
给出一个幺半群 $(S, \cdot)$ 和元素 $u, r \in S$，以及一条直线 $y = \frac{a x + b}{c}$。\\
画出平面中所有坐标为正整数的横线和竖线，维护一个 $f$，初值为单位元 $e$。\\
从原点出发，先向 $y$ 轴正方向走直到到达直线与 $y$ 的交点，然后沿直线走一直走到与 $x = n$ 的交点为止。\\
每当经过一条横线时，执行 $f \gets fu$，经过一条竖线时执行 $f \gets fr$。特别地，在 $y$ 轴上行走时不考虑竖线，同时经过横线和竖线时先执行前者。\\
求最终的 $f$。记为 $\mathrm{euclid}(a, b, c, n, u, r)$。\\
其中 $a, b \ge 0, n, c > 0$。
\end{tcolorbox}

\textbf{做法}：

$$
\mathrm{euclid}(a, b, c, n, u, r) =
\begin{cases}
r ^ n & m = 0 \\
u ^ {\lfloor\frac{b}{c}\rfloor} \cdot \mathrm{euclid}(a \bmod c, b \bmod c, c, n, u, u ^ {\lfloor\frac{a}{c}\rfloor} r) & a \ge c \lor b \ge c \\
r ^ {\lfloor\frac{c - b - 1}{a}\rfloor} u \cdot \mathrm{euclid}(c, (c - b - 1) \bmod a, a, m - 1, r, u) \cdot r ^ {n - \lfloor\frac{c m - b - 1}{a}\rfloor} & \text{otherwise}
\end{cases}
$$

设一次乘法的复杂度为 $O(T)$，则复杂度为 $O(T \log(a + c) \log(a + n + c))$。

\begin{minted}{c++}
// == Preparations ==
struct Node {
    // ...

    Node operator*(const Node &x) const {
        // ...
    }
};
// == Main ==
Node power(Node a, long long b) {
    Node ans = Node(/*幺元*/);
    while (b) {
        if (b & 1) ans = ans * a;
        a = a * a;
        b >>= 1;
    }
    return ans;
}
Node Euclid(int a, int b, int c, long long n, Node r, Node u) {
    long long m = (a * n + b) / c;
    if (!m) return power(r, n);
    if (a >= c || b >= c)
        return power(u, b / c) * Euclid(a % c, b % c, c, n, power(u, a / c) * r, u);
    return power(r, (c - b - 1) / a) * u *
        Euclid(c, (c - b - 1) % a, a, m - 1, u, r) * power(r, n - (c * m - b - 1) / a);
}
\end{minted}

\section{离散对数}

\begin{tcolorbox}
\textbf{定义}：

取有原根的正整数模数 $m$ ，设其一个原根为 $g$ . 对满足 $(a,m)=1$ 的整数 $a$ ，我们知道必存在唯一的整数 $0\leq k<\varphi(m)$ 使得 $ g^k\equiv a\pmod m $。

我们称这个 $k$ 为以 $g$ 为底，模 $m$ 的离散对数，记作 $k=\operatorname{ind}_g a$。
\end{tcolorbox}

离散对数问题即在模 $p$ 意义下求解 $\log_{a}{b}$。这等价于求解离散对数方程

$$a^x \equiv b\pmod p$$

\subsection{BSGS}

运用了根号分治的思想：设块长为 $M$ 且 $x=AM-B$，则有 $a^{AM} \equiv ba^{B} \pmod p$。

固定模数 $p$ 和底数 $a$ 时，预处理时间为 $\mathcal{O}(\frac{p}{M})$，单次询问 $\mathcal{O}(M)$

如果数据组数为 $T$，当 $M$ 取 $\left \lceil \sqrt{\frac{p}{T}} \right \rceil$ 时取到最优复杂度。

以下代码输入要求 $a \perp p,2\le b,n < p<2^{31}$，求得的是模意义下 $\log_{a}{b}$ 的最小的非负整数解。

\begin{minted}{c++}
// == Preparations ==
int a , p , m , t , sz , phi; 
__gnu_pbds::gp_hash_table<int , int>mp;
int Qpow(int x , int p , int mod);
int Phi(int n);
int Inv(int a , int p){return Qpow(a , phi - 1 , p);}
// == Main ==
void Init()
{
    mp.clear();
    phi = Phi(p); 
    sz = sqrt(phi / m) + 1;
    t = ((ll)phi + sz - 1) / sz;
    int as = Qpow(a , sz , p) , aa = 1;
    for(int i = 1 ; i <= t ; i++)
    {
        aa = (ll)aa * as % p;
        if(!mp[aa])mp[aa] = i;
    }
}
int BSGS(int b)
{
    int r = b , ans = p;
    for(int k = 0 ; k <= sz ; k++)
    {
        if(mp.find(r) != mp.end())
            ans = min((ll)ans , (ll)mp[r] * sz - k);
        r = (ll)r * a % p;
    }
    return ans >= p ? -1 : ans;
}
\end{minted}

\subsection{exBSGS}

将方程化为如下形式，其中 $\frac{p}{D} \perp a$。

$$\frac{a^k}{D}\cdot a^{x-k}\equiv\frac{b}{D} \pmod{\frac{p}{D}}$$

然后使用 BSGS 求解，注意特判 $x \le k$ 的情况。

以下代码输入要求 $1\le a,p,b\le10^9$，求得的是模意义下 $\log_{a}{b}$ 的最小的非负整数解。

\begin{minted}{c++}
// == Preparations ==
int a , p , mod , m , t , sz , phi , k , d , inv; 
__gnu_pbds::gp_hash_table<int , int>mp , low;
int Qpow(int x , int p , int mod);
int Phi(int n);
int Inv(int a , int p);
int BSGS(int b);
// == Main ==
void Init()
{
    mp.clear() , low.clear();
    a %= p , mod = p , d = 1 , k = 0; 
    int ad = 1 , ak = 1; 
    for(int g = __gcd(a , p) ; g != 1 ; g = __gcd(a , p))
    {
        ak = (ll)ak * a % mod , k++;
        if(!low[ak])low[ak] = k;
        d *= g , p /= g , ad = ad * ll(a / g) % p;
    }
    phi = Phi(p) , inv = Inv(ad , p);
    sz = sqrt(phi / m) + 1;
    t = ((ll)phi + sz - 1) / sz;
    int as = Qpow(a , sz , p) , aa = 1;
    for(int i = 1 ; i <= t ; i++)
    {
        aa = (ll)aa * as % p;
        if(!mp[aa])mp[aa] = i;
    }
}
int exBSGS(int b)
{
    int bm = b % mod;
    if(mod == 1 || bm == 1)return 0;    
    if(low.find(bm) != low.end())return low[bm];
    if(b % d)return -1;
    b = (ll)(b / d) * inv % p;
    int ans = BSGS(b);
    return ans == -1 ? -1 : ans + k;
}
\end{minted}

\section{原根}

\subsection{阶}

\begin{tcolorbox}
\textbf{定义}：满足同余式 $a^n \equiv 1 \pmod m$ 的最小正整数 $n$ 存在，这个 $n$ 称作 $a$ 模 $m$ 的阶，记作 $\delta_m(a)$ 或 $\operatorname{ord}_m(a)$。
\end{tcolorbox}

\begin{tcolorbox}
\textbf{性质 1}：$a,a^2,\cdots,a^{\delta_m(a)}$ 模 $m$ 两两不同余。
\end{tcolorbox}

\begin{tcolorbox}
\textbf{性质 2}：若 $a^n \equiv 1 \pmod m$ ，则 $\delta_m(a)\mid n$。
\end{tcolorbox}

\begin{tcolorbox}
\textbf{性质 2 推论}：若 $a^p\equiv a^q\pmod m$ ，则有 $p\equiv q\pmod{\delta_m(a)}$。
\end{tcolorbox}

\begin{tcolorbox}
\textbf{性质 3}：

设 $m\in\mathbf{N}^{*}$ ， $a,b\in\mathbf{Z}$ ， $(a,m)=(b,m)=1$ ，则

$$\delta_m(ab)=\delta_m(a)\delta_m(b)$$

的充分必要条件是

$$\left(\delta_m(a), \delta_m(b)\right)=1$$
\end{tcolorbox}

\begin{tcolorbox}
\textbf{性质 4}：

设 $k \in \mathbf{N}$ ， $m$ 为正整数， $a\in\mathbf{Z}$ ， $(a,m)=1$ ，则
$$\delta_m(a^k)=\dfrac{\delta_m(a)}{\left(\delta_m(a),k\right)}$$
\end{tcolorbox}

\subsection{原根}


\begin{tcolorbox}
\textbf{定义}：

设 $m$ 为正整数， $g$ 为整数。 若 $(g,m)=1$ ，且 $\delta_m(g)=\varphi(m)$ ，则称 $g$ 为模 $m$ 的原根。

即 $g$ 满足
$\delta_m(g) = \varphi(m)$。 当 $m$ 是质数时，我们有 $g^i \bmod m,0 < i < m$ 的结果互不相同。
\end{tcolorbox}


\begin{tcolorbox}
\textbf{原根判定定理}：

设 $m \ge 3, (g,m)=1$ ，则 $g$ 是模 $m$ 的原根的充要条件是，对于 $\varphi(m)$ 的每个素因数 $p$ ，都有
$g^{\frac{\varphi(m)}{p}}\not\equiv 1\pmod m$。
\end{tcolorbox}


\begin{tcolorbox}
\textbf{原根个数}：

若一个数 $m$ 有原根，则它原根的个数为 $\varphi(\varphi(m))$。
\end{tcolorbox}


\begin{tcolorbox}
\textbf{原根存在定理}：

一个数 $m$ 存在原根当且仅当 $m=2,4,p^{\alpha},2p^{\alpha}$ ，其中 $p$ 为奇素数， $\alpha\in \mathbf{N}^{*}$。
\end{tcolorbox}

\begin{minted}{c++}
// == Preparations ==
const int N = 1e6 + 5; 
int m , pr[N] , phi[N] , vis[N] , mod; bitset<N>b;
void Sieve(int n = N - 5);
int Qpow(int x , int p , int mod);
// == Main ==
vector<int> PrimitiveRoot(int p , int d)
{
    vector<int>pfac , g; int mod = p;
    if(!vis[p]){return g;} // whether exist primitive root
    for(int i = 1 ; i <= m && pr[i] <= phi[p] ; i++)
        if(phi[p] % pr[i] == 0)pfac.push_back(pr[i]);
    for(int a = 1 ; a < p ; a++) // minimal primitive root is O(n^{0.25+eps})
    {
        if(__gcd(a , p) != 1)continue ;
        bool ok = 1;
        for(int x : pfac)if(Qpow(a , phi[p] / x , mod) == 1){ok = 0; break ;}
        if(!ok)continue ;
        g.push_back(a); break ;
    }
    for(int i = 2 ; i < phi[p] ; i++)
        if(__gcd(i , phi[p]) == 1)g.push_back(Qpow(g[0] , i , mod));
    sort(g.begin() , g.end());
    return g;
}
\end{minted}

\section{剩余}

\begin{tcolorbox}
\textbf{定义}：

令整数 $k\geq 2$ ，整数 $a$ ， $m$ 满足 $(a,m)=1$ ，若存在整数 $x$ 使得
\begin{equation}\tag{1}
x^k\equiv a\pmod m 
\end{equation}

则称 $a$ 为模 $m$ 的 $k$ 次剩余，否则称 $a$ 为模 $m$ 的 $k$ 次非剩余。
\end{tcolorbox}

当整数 $k\geq 2$ ，整数 $a$ ， $m$ 满足 $(a,m)=1$ ，模 $m$ 有原根 $g$ 时，令 $d=(k,\varphi(m))$ ，则：

\begin{tcolorbox}
\textbf{性质 1}：$a$ 为模 $m$ 的 $k$ 次剩余当且仅当 $d\mid \operatorname{ind}_g a$ ，即：$ a^{\frac{\varphi(m)}{d}}\equiv 1\pmod m $
\end{tcolorbox}

\begin{tcolorbox}
\textbf{证明 1}：令 $x=g^y$ ，则方程 $(1)$ 等价于
$ g^{ky}\equiv g^{\operatorname{ind}_g a}\pmod m $

其等价于
$ ky\equiv \operatorname{ind}_g a\pmod{\varphi(m)} $

由同余的性质，我们知道 $y$ 有整数解当且仅当 $d=(k,\varphi(m))\mid \operatorname{ind}_g a$ ，进而
\begin{align} 
\notag a^{\frac{\varphi(m)}{d}}\equiv 1\pmod m
&\notag \iff g^{\frac{\varphi(m)}{d}\operatorname{ind}_g a}\equiv 1\pmod m\\ 
&\notag \iff \varphi(m)\mid\frac{\varphi(m)}{d}\operatorname{ind}_g a\\ 
&\notag \iff \varphi(m)d\mid \varphi(m)\operatorname{ind}_g a\\ 
&\notag \iff d\mid \operatorname{ind}_g a
\end{align}

\end{tcolorbox}

\begin{tcolorbox}
\textbf{性质 2}：方程 $(1)$ 若有解，则模 $m$ 下恰有 $d$ 个解
\end{tcolorbox}

\begin{tcolorbox}
\textbf{性质 3}：模 $m$ 的 $k$ 次剩余类的个数为 $\dfrac{\varphi(m)}{d}$ , 其有形式
    $$ a\equiv g^{di}\pmod m,\qquad \left(0\leq i<\frac{\varphi(m)}{d}\right) $$
\end{tcolorbox}

\subsection{二次剩余}

下面讨论模数是奇素数时的二次剩余问题：

$$x^2 \equiv n \pmod p$$

参考以上结论不难得出，也就说二次剩余的数量恰为 $\frac{p-1}{2}$ ，其他的非 $0$ 数都是非二次剩余，数量也是 $\frac{p-1}{2}$；判别是否为二次剩余只需检查是否 $n^{\frac{p-1}{2}} \equiv 1 \pmod p$ 即可。

\subsection{Cipolla}

通过随机找到一个 $a$ 满足 $a^2 - n$ 是非二次剩余。

定义 $i^2\equiv a^2 - n$，将所有数表示为 $A+Bi$ 的形式，有 $(a+i)^{p+1}\equiv n$。

那么 $(a+i)^{\frac{p-1}{2}}$ 即是一个解，其相反数是另一个解。

\begin{minted}{c++}
// == Preparations ==
typedef long long ll;
int mod , w2;
mt19937 rnd(114514);
struct Complex
{
    int a , b;
    Complex(ll aa = 0 , ll bb = 0):a(aa % mod) , b(bb % mod){}
};
typedef Complex cp;
cp operator + (cp x , cp y){return cp(x.a + y.a , x.b + y.b);}
cp operator - (cp x , cp y){return cp(x.a - y.a , x.b - y.b);}
cp operator * (cp x , cp y)
{
    return cp((ll)x.a * y.a + (ll)w2 * x.b % mod * y.b ,
              (ll)x.b * y.a + (ll)y.b * x.a);
}
cp Qpow(cp x , int p)
{
    cp res = 1;
    while(p)
    {
        if(p & 1)res = res * x;
        x = x * x;
        p >>= 1;
    }
    return res;
}
// == Main ==
int Legendre(int a){return Qpow(a , (mod - 1) / 2).a;}
pair<int , int> Cipolla(int a)
{
    if(a % mod == 0)return {0 , -1};
    if(Legendre(a) == mod - 1)return {-1 , -1};
    int res = 0;
    while(1)
    {
        int r = rnd() % mod;
        w2 = ((ll)r * r - a) % mod;
        if(w2 < 0)w2 += mod;
        if(w2 == 0)
        {
            res = r;
            break ;
        }
        if(Legendre(w2) == mod - 1)
        {
            res = Qpow(cp(r , 1) , (mod + 1) / 2).a;
            break ;
        }
    }
    if(res * 2 >= mod)res = mod - res;
    return {res , mod - res};
}
\end{minted}

\subsection{模素数 N 次剩余}

我们令 $x=g^c$ ， $g$ 是 $p$ 的原根（我们一定可以找到这个 $g$ 和 $c$ ），问题转化为求解 $(g^c)^a \equiv b \pmod p$ . 稍加变换，得到
$$ (g^a)^c \equiv b \pmod p $$

于是就转换成了 BSGS 的基本模型了，可以在 $O(\sqrt p)$ 解出 $c$ ，这样可以得到原方程的一个特解 $x_0\equiv g^c\pmod p$ .



\section{min25 筛}

\subsection{简述}

应用范围：求 $\sum_{i=1}^n f(i)$，其中 $f(i)$ 是积性函数。

需要满足 $f(i)$ 在 $i$ 是质数时的取值是多项式。

时间复杂度：$\Theta(n^{1-\epsilon})$/$\Theta(\frac{n^{\frac{3}{4}}}{\log n})$。

主要想法是将 $f(i)$ 分成三个部分后求和：$i$ 是质数，$i$ 是合数，$i=1$。

先求出 $i$ 是质数部分的和。

为了方便起见，我们将 $f(i)$ 分解成若干 $x^k$ 之和，然后分别计算，这样我们要计算的函数就是\textbf{完全积性函数}了。

考虑埃氏筛的过程，设 $g_{n,k}$ 表示\textbf{筛了 $2 \sim n$ 中的数，用前 $k$ 个质数来筛，没有被筛掉的数的 $f'$ 值之和}。

设 $p_i$ 表示第 $i$ 个质数，考虑求出 $g$ 数组，首先容易发现如果 $p_k^2 > n$，$g_{n,k}=g_{n,k-1}$。

否则我们会筛掉一些数，由于是完全积性函数，容易得出转移：

$$g_{n,k}=g_{n,k-1}-f'(p_k) \times (g_{\lfloor \frac{n}{p_k} \rfloor,k-1}-\sum_{i=1}^{k-1}f'(p_i))$$

注意 $g$ 数组计算的答案是\textbf{“将没被筛掉的数也当成质数”时的答案}。

我们考虑计算答案，设 $s_{n,k}$ 表示\textbf{筛了 $2 \sim n$ 中的数，用前 $k$ 个质数来筛，没有被筛掉的数的 $f$ 值之和。}

现在这个 $s_{n,k}$ 求的就是真正的答案了，最终答案就是 $s_{n,0}+f(1)$。

设 $sp_i=\sum_{j=1}^i p_j$，$D$ 为最大的 $i$ 满足 $p_i^2 \le n$，容易得出转移：

$$s_{n,k}=g_{n,D}-sp_{k}+\sum_{p_i^e \le n \& i>k} f({p_i^e})(s_{\frac{n}{p_i^e},e}+[e>1])$$

然后暴力递归求出 $s$ 就行了，不需要记忆化。

\subsection{求素数个数}

\begin{minted}{c++}
// == Preparations ==
#define int long long
#define id(x) ((x<=lim)?(id1[x]):(id2[n/(x)]))
#define rep(i,j,k) for(int i=j;i<=k;++i)
#define per(i,j,k) for(int i=j;i>=k;--i)
int const N=1e6+10;
int lim,prime[N],w[N],g[N],id1[N],id2[N],m;bool vis[N];
// == Main ==
inline void init(){
    //预处理根号以内质数
    rep(i,2,lim){
        if (!vis[i]) prime[++prime[0]]=i;
        for (int j=1;j<=prime[0] && i*prime[j]<=lim;++j){
            vis[i*prime[j]]=1;
            if (i%prime[j]==0) break;
        }
    }
}
inline void solve(){
    int n;cin>>n,lim=sqrtl(n),init();
    for (int l=1,r;l<=n;l=r+1){
        r=n/(n/l),w[++m]=n/l,g[m]=w[m]-1;
        if (w[m]<=lim) id1[w[m]]=m;
        else id2[n/w[m]]=m;
        //预处理 g 数组第一维有可能的取值
    }
    rep(i,1,prime[0])
        for (int j=1;j<=m && prime[i]*prime[i]<=w[j];++j)
            g[j]-=g[id(w[j]/prime[i])]-(i-1);
    cout<<g[id(n)]<<'\n';
}
\end{minted}

\subsection{完整模板}

\begin{minted}{c++}
// == Preparations ==
#define int long long
#define rep(i,j,k) for(int i=j;i<=k;++i)
#define per(i,j,k) for(int i=j;i>=k;--i)
#define id(x) ((x<=lim)?(id1[x]):(id2[n/(x)]))
#define add(x,y) (x=((x+y>=mod)?(x+y-mod):(x+y)))
// == Main ==
//例：求 f(p^k)=p^k(p^k-1) 的前缀和
//发现 f(p^k)=p^{2k}-p^k
int const mod=1e9+7;
int n;
inline int f(int op,int x){
    x%=mod;
    if (!op) return x;
    return x*x%mod;
}
inline int F(int x){
    x%=mod;
    return x*(x-1)%mod;
}
inline int smf(int op,int x){
    x%=mod;
    if (!op) return (x*(x+1)/2)%mod;
    return x*(x+1)%mod*(2*x%mod+1)%mod*166666668%mod;
}
//smf -> f 前缀和
int const N=1e6+10;
int lim,prime[N],sp1[N],sp2[N],w[N],g1[N],g2[N],id1[N],id2[N],m;bool vis[N];
inline void init(){
    //预处理根号以内质数
    rep(i,2,lim){
        if (!vis[i]) prime[++prime[0]]=i;
        for (int j=1;j<=prime[0] && i*prime[j]<=lim;++j){
            vis[i*prime[j]]=1;
            if (i%prime[j]==0) break;
        }
    }
    rep(i,1,prime[0]){
        sp1[i]=sp1[i-1],add(sp1[i],f(0,prime[i]));
        sp2[i]=sp2[i-1],add(sp2[i],f(1,prime[i]));
    }
}
inline int S(int x,int y){
    if (prime[y]>=x) return 0;
    int ans=((g2[id(x)]+mod-g1[id(x)])%mod+mod-(sp2[y]-sp1[y]+mod)%mod)%mod;
    for (int i=y+1;i<=prime[0] && prime[i]*prime[i]<=x;++i)
        for (int e=1,gg=prime[i];gg<=x;++e,gg*=prime[i])
            add(ans,F(gg)*(S(x/gg,i)+(e>1))%mod);
    return ans;
}
inline void solve(){
    cin>>n,lim=sqrtl(n),init();
    for (int l=1,r;l<=n;l=r+1){
        r=n/(n/l),w[++m]=n/l;
        g1[m]=(smf(0,w[m])+mod-1)%mod;
        g2[m]=(smf(1,w[m])+mod-1)%mod;
        if (w[m]<=lim) id1[w[m]]=m;
        else id2[n/w[m]]=m;
        //预处理 g 数组第一维有可能的取值
    }
    rep(i,1,prime[0])
        for (int j=1;j<=m && prime[i]*prime[i]<=w[j];++j)
            add(g1[j],mod-f(0,prime[i])*(g1[id(w[j]/prime[i])]+mod-sp1[i-1])%mod),
            add(g2[j],mod-f(1,prime[i])*(g2[id(w[j]/prime[i])]+mod-sp2[i-1])%mod);
    cout<<(S(n,0)+1)%mod<<'\n';
}
\end{minted}


%
%% Chapter ? 线性代数
\chapter{线性代数}

\section{线性空间}

\begin{tcolorbox}
    \textbf{线性空间}：给定域 $F$，$F$ 上的线性空间 $V$ 是一个集合，其上定义了两种二元运算：
    \begin{itemize}
        \item 加法 $+: V\times V\to V$；
        \item 数乘 $\cdot: F\times V\to V$。
    \end{itemize}
    且满足：
    \begin{multicols}{2}
        \begin{itemize}
            \item $\forall \mathbf{u},\mathbf{v},\mathbf{w}\in V,\mathbf{u}+(\mathbf{v}+\mathbf{w})=(\mathbf{u}+\mathbf{v})+\mathbf{w}$
            \item $\forall \mathbf{u},\mathbf{v}\in V,\mathbf{u}+\mathbf{v}=\mathbf{v}+\mathbf{u}$
            \item $\exists \mathbf{0}\in V,\forall \mathbf{v}\in V,\mathbf{v}+\mathbf{0}=\mathbf{0}+\mathbf{v}=\mathbf{v}$
            \item $\forall \mathbf{v}\in V,\exists -\mathbf{v}\in V,\mathbf{v}+(-\mathbf{v})=\mathbf{0}$
            \item $\exists 1\in F,\forall \mathbf{v}\in V,1\mathbf{v}=\mathbf{v}$
            \item $\forall a,b\in F,\mathbf{v}\in V,a(b\mathbf{v})=(ab)\mathbf{v}$
            \item $\forall a\in F,\mathbf{u},\mathbf{v}\in V,a(\mathbf{u}+\mathbf{v})=a\mathbf{u}+a\mathbf{v}$
            \item $\forall a,b\in F,\mathbf{v}\in V,(a+b)\mathbf{v}=a\mathbf{v}+b\mathbf{v}$
        \end{itemize}
    \end{multicols}
\end{tcolorbox}

\subsection{正交子空间与正交补}

\begin{tcolorbox}
    \textbf{正交子空间}：若一个内积空间的子空间 $A,B$ 满足一者中的每个向量都与另一者正交（即内积为 $0$），那么它们互为正交子空间。
\end{tcolorbox}

其中内积空间是增添了一个额外的结构（即内积）的向量空间，这里略过其定义，可以狭义地理解为欧几里得空间 $\mathbb{R}^n$，内积定义为点积。

\begin{tcolorbox}
    \textbf{正交补}：对于一个内积空间的子空间 $A$，若另一个子空间 $B$ 包含所有与子空间 $A$ 正交的向量，则称 $B$ 为 $A$ 的正交补。
\end{tcolorbox}

\section{矩阵}

\subsection{矩阵乘法的另一种理解}

\begin{itemize}
    \item $A$ 右乘矩阵 $B$，是对 $A$ 的列向量进行线性组合，组合系数记在 $B$ 的每个列向量中。
    \item $B$ 左乘矩阵 $A$，是对 $B$ 的行向量进行线性组合，组合系数记在 $A$ 的每个行向量中。
\end{itemize}

\subsection{基本子空间}

\begin{tcolorbox}
    \textbf{行空间}：矩阵 $A$ 的行向量组的生成子空间称为行空间，记为 $C(A^{\mathrm{T}})$。

    \textbf{列空间}：矩阵 $A$ 的列向量组的生成子空间称为列空间，记为 $C(A)$。
\end{tcolorbox}
\begin{tcolorbox}
    根据矩阵的秩的定义，有：对于任意矩阵 $A$，有 $\dim C(A^{\mathrm{T}})=\dim C(A)=\operatorname{rank}(A)$。
\end{tcolorbox}

\begin{tcolorbox}
    \textbf{零空间}：矩阵 $A$ 的零空间为 $Ax=0$ 的所有解构成的空间，记为 $N(A)$。

    \textbf{左零空间}：矩阵 $A$ 的左零空间为 $A^{\mathrm{T}}x=0$（或表述为 $x^{\mathrm{T}}A=0$） 的所有解构成的空间，记为 $N(A^{\mathrm{T}})$。
\end{tcolorbox}

\begin{tcolorbox}
    有：对于任意 $n\times m$ 的矩阵 $A$，有 $\dim N(A)=m-\operatorname{rank}(A),\dim N(A^{\mathrm{T}})=n-\operatorname{rank}(A)$。
\end{tcolorbox}

\begin{tcolorbox}
    \textbf{四个基本子空间的正交性}：根据四个子空间的定义，对于任意矩阵 $A$ 有：
    \begin{itemize}
        \item $N(A)$ 与 $C(A^{\mathrm{T}})$ 互为正交子空间，且互为正交补。
        \item $N(A^{\mathrm{T}})$ 与 $C(A)$ 互为正交子空间，且互为正交补。
    \end{itemize}

\end{tcolorbox}

\subsection{高斯消元}

\begin{minted}{c++}
// == Preparations ==
double a[M][N];
int Cmp(double x , double y = 0){return abs(x - y) < 1e-8 ? 0 : (x < y ? -1 : 1);}
// == Main ==
int Gauss()
{
    int tot = 0;
    for(int i = 1 ; i <= n ; i++)
    {
        for(int j = tot + 1 ; j <= m ; j++)
        {
            if(Cmp(a[j][i]))
            {
                for(int k = 1 ; k <= n + 1 ; k++)
                    swap(a[tot + 1][k] , a[j][k]);
                break ;
            }
        }
        if(!Cmp(a[tot + 1][i]))continue ;
        for(int j = tot + 2 ; j <= m ; j++)
        {
            double r = a[j][i] / a[tot + 1][i];
            for(int k = 1 ; k <= n + 1 ; k++)
                a[j][k] -= r * a[tot + 1][k];   
        }
        tot++;
    }
    for(int i = tot + 1 ; i <= m ; i++)
    {
        if(Cmp(a[i][n + 1]))
        {
            return -1;// No solutions
        }
    }
    if(tot == n)
    {
        for(int i = n ; i >= 1 ; i--)
        {
            a[i][n + 1] /= a[i][i] , a[i][i] = 1;   
            for(int j = i - 1 ; j >= 1 ; j--)
                a[j][n + 1] -= a[j][i] * a[i][n + 1] , a[j][i] = 0;
        }
        return 0;// One solution
    }
    return 1;// Many solutions
}
\end{minted}

\subsection{逆矩阵}

\begin{tcolorbox}
    \textbf{逆矩阵}：设 $A$ 是一个 $n$ 阶方阵，若存在 $n$ 阶方阵 $B$ 使得 $AB=BA=I_n$，其中 $I_n$ 为 $n$ 阶单位矩阵，则称 $A$ 是可逆的，是非奇异矩阵，$B$ 是 $A$ 的逆矩阵，记为 $A^{-1}$。否则称 $A$ 不可逆，是奇异矩阵。
\end{tcolorbox}

逆矩阵可以用来快速求解线性方程组。设有 $Ax=b$，则有 $x=A^{-1}b$，对于不同的 $b$，每次只需要 $\mathcal{O}(n^2)$ 即可计算一组解。这样做的前提是 $A$ 可逆。

\begin{minted}{c++}
// == Preparations ==
const int MOD = 1e9 + 7;
int Inv(int x);
// == Main ==
bool Inv(int n , int a[N][N * 2])
{
    for(int i = 1 ; i <= n ; i++)
        a[i][i + n] = 1;
    for(int i = 1 ; i <= n ; i++)
    {
        for(int j = i ; j <= n ; j++)
        {
            if(a[j][i] != 0)
            {
                for(int k = 1 ; k <= n * 2 ; k++)
                    swap(a[j][k] , a[i][k]);
                break ;
            }
        }
        if(!a[i][i])return 0;// No Solution
        int inv = Inv(a[i][i]);
        for(int j = 1 ; j <= n ; j++)
        {
            if(j == i)continue ;
            int r = (ll)a[j][i] * inv % MOD;
            for(int k = 1 ; k <= n * 2 ; k++)
                a[j][k] = (a[j][k] - (ll)a[i][k] * r) % MOD;
        }
        for(int j = 1 ; j <= n * 2 ; j++)
            a[i][j] = (ll)a[i][j] * inv % MOD;
    }
    for(int i = 1 ; i <= n ; i++)
        for(int j = 1 ; j <= n ; j++)
            if(a[i][j + n] < 0)a[i][j + n] += MOD;
    return 1;
}
\end{minted}

\section{行列式}

\begin{tcolorbox}
    \textbf{行列式}：一个 $n$ 阶方阵 $A$ 的行列式定义为
    $$\sum_{\sigma\in S_n} (-1)^{\operatorname{N}(\sigma)}\prod_{i=1}^{n} A_{i,\sigma_i}$$
    其中 $S_n$ 是所有 $n$ 阶排列的全体，$\operatorname{N}(\sigma)$ 表示 $\sigma$ 的逆序对数。

    我们用 $\det(A)$ 或 $|A|$ 表示 $A$ 的行列式。
\end{tcolorbox}

\subsection{行列式的基本性质}

\begin{itemize}
    \item 某一行（列）有公因子 $k$，则可以提出 $k$。
    \item 两行（列）互换，改变行列式正负符号：有向面积的方向改变。
    \item 在行列式中，有两行（列）对应成比例或相同，则此行列式的值为 $0$：有两个重叠的向量，有向面积为零。
    \item 在行列式中，某一行（列）的每个元素是两数之和，则此行列式可拆分为两个相加的行列式：一行只能选一个数，这个数拆成和的左 / 右部分。
    \item 将一行(列)的 $k$ 倍加进另一行(列)里，行列式的值不变：可用上两条性质证明。
    \item 将行列式的行列互换，行列式的值不变。
\end{itemize}

\subsection{余子式与拉普拉斯展开}

\begin{tcolorbox}
    对于方阵 $A$，$I=\{i_1,i_2,\ldots,i_k\}\subseteq R$，$J=\{j_1,j_2,\ldots,j_k\}\subseteq C$，定义：

    \begin{itemize}
        \item $|A[I,J]|$ 称为子式，即选出一个子方阵求行列式。若 $I=J$（即选出的行、列集合相同），则也可称为主子式。当 $I=J=\{1,2,\ldots,k\}$ 时，也称为 $k$ 阶顺序主子式。
        \item $|A[R\setminus I,C\setminus J]|$ 称为余子式。
        \item $\operatorname{sgn}(I,J)|A[R\setminus I,C\setminus J]|$ 称为余子式。
    \end{itemize}
\end{tcolorbox}

其中 $\operatorname{sgn}(I,J)=(-1)^{i_1+\cdots+i_k+j_1+\cdots+j_k}$，下同。

矩阵的秩等于最高阶非零子式的阶数。这是因为如果子矩阵存在线性相关的两行（列），则行列式为零。

\begin{tcolorbox}
    \textbf{拉普拉斯展开}：对于 $n$ 阶方阵 $A$，对于任意 $1\le i\le n$，有
    $$|A|=\sum_{j=1}^{n} (-1)^{i+j}|A[R\setminus\{i\},C\setminus \{j\}]|A_{i,j}$$
\end{tcolorbox}
\begin{itemize}
    \item 对列也有同样的性质。
    \item 根据行列式的定义即可证明。
\end{itemize}

\begin{tcolorbox}
    \textbf{拉普拉斯定理}：对于方阵 $A$，以及任意 $I\subseteq R$，有
    $$|A|=\sum_{J\subseteq C,|J|=|I|} \operatorname{sgn}(I,J)|A[I,J]|\cdot|A[R\setminus I,C\setminus J]|$$
\end{tcolorbox}

\begin{minted}{c++}
// == Preparations ==
typedef long long ll;
const int N = 605;
// == Main ==
int Det(int n , int mod , int a[N][N])
{
    int m = 0 , f = 1;
    auto Swap = [&](int i , int j)
    {
        if(i == j)return ;
        for(int k = 1 ; k <= n ; k++)
            swap(a[i][k] , a[j][k]);
        f = -f;
    };
    for(int i = 1 ; i <= n ; i++)
    {
        m++;
        for(int j = m ; j <= n ; j++)
            if(a[j][i]){Swap(m , j); break ;};
        if(!a[m][i])return 0;
        for(int j = m + 1 ; j <= n ; j++)
        {
            while(a[j][i] && a[m][i])
            {
                ll r = a[j][i] / a[m][i] % mod;
                for(int k = 1 ; k <= n ; k++)
                    a[j][k] = (a[j][k] - r * a[m][k]) % mod;
                Swap(m , j);
            }
            if(!a[m][i])Swap(m , j);
        }
    }
    ll ans = 1;
    for(int i = 1 ; i <= n ; i++)
        (ans *= a[i][i]) %= mod;
    return (ans * f % mod + mod) % mod;
}
\end{minted}
%
%% Chapter ? 图论
\chapter{图论}

\section{k 短路}

复杂度为 $O((n + m)\log n + k\log k)$。

\begin{minted}{c++}
// == Preparations ==
int ontree[200005];
struct graph {
    int tot, hd[5005];
    int nxt[200005], to[200005];
    long long dt[200005];

    void add(int u, int v, long long w) {
        nxt[++tot] = hd[u];
        hd[u] = tot;
        to[tot] = v;
        dt[tot] = w;
        return ;
    }
} g;
long long dis[5005];
struct node {
    int id;
    long long val;

    node() = default;
    node(int _id, long long _val): id(_id), val(_val) {}
    bool operator<(const node &x) const {return val > x.val;}
};
priority_queue<node> q;
int vis[5005];

// 以下左偏树
struct HeapNode {
    long long val;
    int to, dist;
    HeapNode *ls, *rs;

    HeapNode() = default;
    HeapNode(long long _val, int _to): val(_val), to(_to), dist(1), ls(nullptr), rs(nullptr) {}
};
struct Heap {
    HeapNode *root[5005];

    HeapNode *merge(HeapNode *u, HeapNode *v) {
        if (!u) return v;
        if (!v) return u;
        if (u->val > v->val) swap(u, v);
        HeapNode *p = new HeapNode(*u);
        p->rs = merge(u->rs, v);
        if (!p->ls || p->ls->dist < p->rs->dist) swap(p->ls, p->rs);
        if (p->rs) p->dist = p->rs->dist + 1;
        else p->dist = 1;
        return p;
    }
} h;

struct Node {
    HeapNode *id;
    long long val;

    Node() = default;
    Node(HeapNode *_id, long long _val): id(_id), val(_val) {}
    bool operator<(const Node &x) const {return val > x.val;}
};
priority_queue<Node> Q;
// == Main ==
void dfs(int now) {
    vis[now] = 1;
    for (int i = g.hd[now]; i; i = g.nxt[i])
        if (!vis[g.to[i]] && dis[g.to[i]] == dis[now] + g.dt[i]) ontree[i] = 1, dfs(g.to[i]);
    return ;
}
void dfs2(int now) {
    for (int i = g.hd[now]; i; i = g.nxt[i])
        if (ontree[i]) h.root[g.to[i]] = h.merge(h.root[g.to[i]], h.root[now]), dfs2(g.to[i]);
    return;
}

memset(dis, 0x3f, sizeof(dis));
q.emplace(n, dis[n] = 0);
while (!q.empty()) {
    int now = q.top().id;
    long long tmp = q.top().val;
    q.pop();
    if (tmp != dis[now]) continue;
    for (int i = g.hd[now]; i; i = g.nxt[i])
        if (dis[g.to[i]] > dis[now] + g.dt[i]) q.emplace(g.to[i], dis[g.to[i]] = dis[now] + g.dt[i]);
}
dfs(n);
for (int i = 1; i <= n; i++)
    for (int j = g.hd[i]; j; j = g.nxt[j])
        if (!ontree[j] && g.to[j] != n)
            h.root[g.to[j]] = h.merge(h.root[g.to[j]], new HeapNode(dis[i] + g.dt[j] - dis[g.to[j]], i));
dfs2(n);
if (h.root[1]) Q.emplace(h.root[1], dis[1] + h.root[1]->val);
while (!Q.empty()) { // 每次取出来一条路径
    HeapNode *now = Q.top().id;
    long long d = Q.top().val;
    Q.pop();
    if (now->ls) Q.emplace(now->ls, d - now->val + now->ls->val);
    if (now->rs) Q.emplace(now->rs, d - now->val + now->rs->val);
    HeapNode *tmp = h.root[now->to];
    if (tmp) Q.emplace(tmp, d + tmp->val);
}
\end{minted}
%
%% Chapter ? 多项式
\chapter{多项式}

\section{拉格朗日插值法}

$$
f(k) = \sum\limits_{i = 0}^n y_i \prod\limits_{j \neq i} \dfrac{k - x_j}{x_i - x_j}
$$

取值连续时的优化：

当 $x_i = i$ 时可以优化成 $O(n)$。

$$
\begin{aligned}
f(k) &= \sum\limits_{i = 0}^n y_i \prod\limits_{j \neq i} \dfrac{k - x_j}{x_i - x_j} \\
&= \sum\limits_{i = 0}^n y_i \prod\limits_{j \neq i} \dfrac{k - j}{i - j}
\end{aligned}
$$

令 $pre_i = \prod\limits_{j = 0}^i (k - j), suf_i = \prod\limits_{j = i}^n (k - j)$：

$$
f(k) = \sum\limits_{i = 0}^n y_i \dfrac{pre_{i - 1} \times suf_{i + 1}}{i! \times (n - i)!} \times (-1)^{n - i}
$$

重心拉格朗日插值：

$$
\begin{aligned}
f(k) &= \sum\limits_{i = 0}^n y_i \prod\limits_{i \neq j} \dfrac{k - x_j}{x_i - x_j} \\
&= \sum\limits_{i = 0}^n y_i \dfrac{\prod\limits_{j \neq i} k - x_j}{\prod\limits_{j \neq i} x_i - x_j} \\
&= \sum\limits_{i = 0}^n \dfrac{y_i}{k - x_i} \dfrac{\prod\limits_{j = 0}^n k - x_j}{\prod\limits_{j \neq i} x_i - x_j} \\
\end{aligned}
$$

设 $g(k) = \prod\limits_{i = 0}^n k - x_i, t_i = \dfrac{y_i}{\prod\limits_{j \neq i} x_i - x_j}$。

则

$$
\begin{aligned}
f(k) &= \sum\limits_{i = 0}^n \dfrac{g(k)t_i}{k - x_i} \\
&= g(k) \sum\limits_{i = 0}^n \dfrac{t_i}{k - x_i}
\end{aligned}
$$

于是每次插入只需要更新之前的所有 $t_i$ 并计算当前点的 $t_i$ 即可，时间复杂度为 $O(n)$。

由于 $g(k)$ 和 $\sum\limits_{i = 0}^k \dfrac{t_i}{k - x_i}$ 都可以 $O(n)$ 求得，故询问复杂度也是 $O(n)$。


\section{牛顿迭代}

用于解决下列问题：

\begin{tcolorbox}
已知函数 $G$ 且 $G(F(x)) = 0$，求多项式 $F$（${}\bmod x^n$）。
\end{tcolorbox}

\textbf{结论}：

$$
F(x) = F_*(x) - \frac{G(F_*(x))}{G'(F_*(x))} \pmod{x^n}
$$

其中 $F_*(x)$ 为做到 $x^{n / 2}$ 时的答案。

\section{FFT}

\begin{minted}{c++}
// == Preparations ==
struct complex {
    double a, b;

    complex() = default;
    complex(double _a, double _b): a(_a), b(_b) {}
    complex operator+(const complex &x) const {return complex(a + x.a, b + x.b);}
    complex operator-(const complex &x) const {return complex(a - x.a, b - x.b);}
    complex operator*(const complex &x) const {return complex(a * x.a - b * x.b, a * x.b + b * x.a);}
    complex operator/(const complex &x) const {
        double t = b * b + x.b * x.b;
        return complex((a * x.a + b * x.b) / t, (b * x.a - a * x.b) / t);
    }
    complex &operator+=(const complex &x) {return *this = *this + x;}
    complex &operator-=(const complex &x) {return *this = *this - x;}
    complex &operator*=(const complex &x) {return *this = *this * x;}
    complex &operator/=(const complex &x) {return *this = *this / x;}
};
// == Main ==
void FFT(vector<complex> &f, int flag) const {
    int n = f.size();
    vector<int> swp(n);
    for (int i = 0; i < n; i++) {
        swp[i] = swp[i >> 1] >> 1 | ((i & 1) * (n >> 1));
        if (i < swp[i]) std::swap(f[i], f[swp[i]]);
    }
    for (int mid = 1; mid < n; mid <<= 1) {
        complex w1(cos(pi / mid), flag * sin(pi / mid));
        for (int i = 0; i < n; i += mid << 1) {
            complex w(1, 0);
            for (int j = 0; j < mid; j++, w *= w1) {
                complex x = f[i + j], y = w * f[i + mid + j];
                f[i + j] = x + y, f[i + mid + j] = x - y;
            }
        }
    }
    return;
}
\end{minted}

\section{常用 NTT 模数及其原根}

\begin{center}
    \begin{tabular}{c|c|c}
        % \specialrule{2pt}{0em}{0em}
        模数 & 原根 & 分解 \\
        \hline
        % \specialrule{1pt}{0em}{0em}
        $167772161$ & $3$ & $5 \times 2^{25} + 1$ \\
        \hline
        $469762049$ & $3$ & $7 \times 2^{26} + 1$ \\
        \hline
        $998244353$ & $3$ & $119 \times 2^{23} + 1$ \\
        \hline
        $1004535809$ & $3$ & $479 \times 2^{21} + 1$ \\
        \hline
        $2013265921$ & $31$ & $15 \times 2^{27} + 1$ \\
        \hline
        $2281701377$ & $3$ & $17 \times 2^{27} + 1$ \\
        % \specialrule{2pt}{0em}{0em}
    \end{tabular}
\end{center}

\section{多项式模板}

\begin{minted}{c++}
// == Preparations ==
#include <vector>
// == Main ==
namespace Poly {
    const int mod = 998244353, G = 3, invG = 332748118;

    inline int power(int a, int b) {
        int ans = 1;
        while (b) {
            if (b & 1) ans = (long long)ans * a % mod;
            a = (long long)a * a % mod;
            b >>= 1;
        }
        return ans % mod;
    }

    struct poly: vector<int> {
        poly(initializer_list<int> &&arg): vector<int>(arg) {}
        template<typename... argT>
        poly(argT &&...args): vector<int>(forward<argT>(args)...) {}

        poly operator+(const poly &b) const {
            const poly &a = *this;
            poly ans(max(a.size(), b.size()));
            for (int i = 0; i < (int)ans.size(); i++)
                ans[i] = ((i < (int)a.size() ? a[i] : 0) + (i < (int)b.size() ? b[i] : 0)) % mod;
            return ans;
        }
        poly operator+=(const poly &b) {return *this = *this + b;}
        poly operator-(const poly &b) const {
            const poly &a = *this;
            poly ans(max(a.size(), b.size()));
            for (int i = 0; i < (int)ans.size(); i++)
                ans[i] = ((i < (int)a.size() ? a[i] : 0) - (i < (int)b.size() ? b[i] : 0) + mod) % mod;
            return ans;
        }
        poly operator-=(const poly &b) {return *this = *this - b;}
        void NTT(poly &g, int flag) const {
            int n = g.size();
            vector<unsigned long long> f(g.begin(), g.end());
            vector<int> swp(n);
            for (int i = 0; i < n; i++) {
                swp[i] = swp[i >> 1] >> 1 | ((i & 1) * (n >> 1));
                if (i < swp[i]) std::swap(f[i], f[swp[i]]);
            }
            for (int mid = 1; mid < n; mid <<= 1) {
                int w1 = power(flag ? G : invG, (mod - 1) / mid / 2);
                vector<int> w(mid);
                w[0] = 1;
                for (int i = 1; i < mid; i++) w[i] = (long long)w[i - 1] * w1 % mod;
                for (int i = 0; i < n; i += mid << 1)
                    for (int j = 0; j < mid; j++) {
                        int t = (long long)w[j] * f[i + mid + j] % mod;
                        f[i + mid + j] = f[i + j] - t + mod;
                        f[i + j] += t;
                    }
                if (mid == 1 << 10)
                    for (int i = 0; i < n; i++) f[i] %= mod;
            }
            int inv = flag ? 1 : power(n, mod - 2);
            for (int i = 0; i < n; i++) g[i] = f[i] % mod * inv % mod;
            return;
        }

        // 下面是基于转置原理的 NTT，相对朴素版本效率更高。
        void NTT(poly &g, int flag) const {
            int n = g.size();
            vector<int> f(g.begin(), g.end());
            if (flag) {
                for (int mid = n >> 1; mid >= 1; mid >>= 1) {
                    int w1 = power(G, (mod - 1) / mid / 2);
                    vector<int> w(mid);
                    w[0] = 1;
                    for (int i = 1; i < mid; i++) w[i] = (long long)w[i - 1] * w1 % mod;
                    for (int i = 0; i < n; i += mid << 1)
                        for (int j = 0; j < mid; j++) {
                            int t = (long long)(f[i + j] - f[i + mid + j] + mod) * w[j] % mod;
                            f[i + j] = f[i + j] + f[i + mid + j] >= mod ?
                                f[i + j] + f[i + mid + j] - mod : f[i + j] + f[i + mid + j];
                            f[i + mid + j] = t;
                        }
                }
                for (int i = 0; i < n; i++) g[i] = f[i];
            } else {
                for (int mid = 1; mid < n; mid <<= 1) {
                    int w1 = power(invG, (mod - 1) / mid / 2);
                    vector<int> w(mid);
                    w[0] = 1;
                    for (int i = 1; i < mid; i++) w[i] = (long long)w[i - 1] * w1 % mod;
                    for (int i = 0; i < n; i += mid << 1)
                        for (int j = 0; j < mid; j++) {
                            int t = (long long)w[j] * f[i + mid + j] % mod;
                            f[i + mid + j] = f[i + j] - t < 0 ? f[i + j] - t + mod : f[i + j] - t;
                            f[i + j] = f[i + j] + t >= mod ? f[i + j] + t - mod : f[i + j] + t;
                        }
                }
                int inv = power(n, mod - 2);
                for (int i = 0; i < n; i++) g[i] = (long long)f[i] * inv % mod;
            }
            return;
        }

        poly operator*(poly b) const {
            poly a(*this);
            int n = 1, len = (int)(a.size() + b.size()) - 1;
            while (n < len) n <<= 1;
            a.resize(n), b.resize(n);
            NTT(a, 1), NTT(b, 1);
            poly c(n);
            for (int i = 0; i < n; i++) c[i] = (long long)a[i] * b[i] % mod;
            NTT(c, 0);
            c.resize(len);
            return c;
        }
        poly operator*=(const poly &b) {return *this = *this * b;}
        poly inv() const {
            poly f = *this, g;
            g.push_back(power(f[0], mod - 2));
            int n = 1;
            while (n < (int)f.size()) n <<= 1;
            f.resize(n << 1);
            for (int len = 2; len <= n; len <<= 1) {
                poly tmp(len), ff(len << 1);
                for (int i = 0; i < len >> 1; i++) tmp[i] = g[i] * 2 % mod;
                for (int i = 0; i < len; i++) ff[i] = f[i];
                g.resize(len << 1);
                NTT(g, 1), NTT(ff, 1);
                for (int i = 0; i < len << 1; i++) g[i] = (long long)g[i] * g[i] % mod * ff[i] % mod;
                NTT(g, 0);
                g.resize(len);
                for (int i = 0; i < len; i++) g[i] = (tmp[i] - g[i] + mod) % mod;
            }
            g.resize(size());
            return g;
        }
        poly sqrt() const { // need F[0] = 1.
            poly f = *this, g;
            g.push_back(1);
            int n = 1;
            while (n < (int)f.size()) n <<= 1;
            f.resize(n << 1);
            for (int len = 2; len <= n; len <<= 1) {
                poly tmp(len), ff(len << 1);
                for (int i = 0; i < len >> 1; i++) tmp[i] = g[i] * 2 % mod;
                for (int i = 0; i < len; i++) ff[i] = f[i];
                g.resize(len << 1);
                NTT(g, 1);
                for (int i = 0; i < len << 1; i++) g[i] = (long long)g[i] * g[i] % mod;
                NTT(g, 0);
                g += ff;
                g *= tmp.inv();
                g.resize(len);
            }
            g.resize(size());
            return g;
        }
        poly derivative() const {
            poly f(*this);
            for (int i = 1; i < (int)f.size(); i++) f[i - 1] = (long long)f[i] * i % mod;
            f.pop_back();
            return f;
        }
        poly integral() const {
            poly f(*this);
            f.push_back(0);
            for (int i = f.size() - 1; i >= 1; i--) f[i] = (long long)f[i - 1] * power(i, mod - 2) % mod;
            f[0] = 0;
            return f;
        }
        poly ln() const {
            poly f((derivative() * inv()).integral());
            f.resize(size()); 
            return f;
        }
        poly exp() const { // 需要满足 F[0] = 0
            poly f(*this), g;
            g.push_back(1);
            int n = 1;
            while (n < (int)size()) n <<= 1;
            f.resize(n);
            for (int len = 2; len <= n; len <<= 1) {
                poly tmp(g);
                g.resize(len);
                g = g.ln();
                for (int i = 0; i < len; i++) g[i] = (f[i] - g[i] + mod) % mod;
                g[0] = (g[0] + 1) % mod;
                g *= tmp;
                g.resize(len);
            }
            g.resize(size());
            return g;
        }
    };

    inline poly power(poly f, int b) { // 需要满足 F[0] = 1
        f = f.ln();
        for (int i = 0; i < (int)f.size(); i++) f[i] = (long long)f[i] * b % mod;
        f = f.exp();
        return f;
    }
    // 不要求 F[0] = 1 的多项式快速幂
    /*
    b1 为指数 mod m，b2 为指数 mod (phi(m) = mod - 1)
    如果 b1 传入之前被取模了记得特判 a[0] = 0 且指数 > n 的情况，此时多项式每一位系数都是 0
    */
    poly power(poly f, int b1, int b2 = -1) {
        if (b2 == -1) b2 = b1;
        int n = f.size(), p = 0;
        reverse(f.begin(), f.end());
        while (!f.empty() && !f.back()) f.pop_back(), p++;
        if (f.empty() || (long long)p * b1 >= n) return poly(n);
        int v = f.back();
        int inv = power(v, mod - 2);
        for (int &i : f) i = (long long)i * inv % mod;
        reverse(f.begin(), f.end());
        f = f.ln();
        for (int &i : f) i = (long long)i * b1 % mod;
        f = f.exp();
        reverse(f.begin(), f.end());
        for (int i = 1; i <= p * b1; i++) f.push_back(0);
        reverse(f.begin(), f.end());
        f.resize(n);
        v = power(v, b2);
        for (int &i : f) i = (long long)i * v % mod;
        return f;
    }
    poly div(poly f, poly g) {
        int n = f.size() - 1, m = g.size() - 1;
        reverse(f.begin(), f.end()), reverse(g.begin(), g.end());
        g.resize(n - m + 1);
        poly ans = f * g.inv();
        ans.resize(n - m + 1);
        reverse(ans.begin(), ans.end());
        return ans;
    }
}
\end{minted}

\section{FWT}

\begin{minted}{c++}
void FWTor(int a[], int op) {
    for (int l = 2; l <= n; l <<= 1)
        for (int i = 0; i < n; i += l)
            for (int j = 0; j < l / 2; j++)
                a[i + j + l / 2] = (a[i + j + l / 2] + (long long)a[i + j] * op % mod + mod) % mod;
    return;
}
void FWTand(int a[], int op) {
    for (int l = 2; l <= n; l <<= 1)
        for (int i = 0; i < n; i += l)
            for (int j = 0; j < l / 2; j++)
                a[i + j] = (a[i + j] + (long long)a[i + j + l / 2] * op % mod + mod) % mod;
    return;
}
void FWTxor(int a[], int op) {
    for (int l = 2; l <= n; l <<= 1)
        for (int i = 0; i < n; i += l)
            for (int j = 0; j < l / 2; j++) {
                int p = (a[i + j] + a[i + j + l / 2]) % mod, d = (a[i + j] - a[i + j + l / 2] + mod) % mod;
                a[i + j] = (long long)p * op % mod;
                a[i + j + l / 2] = (long long)d * op % mod;
            }
    return;
}
\end{minted}
%
%% Chapter ? 计算几何
\chapter{计算几何}

\section{计算几何模板}

\begin{minted}{c++}
// == Preparations ==
#include <cmath>
#include <vector>
#include <algorithm>
// == Main ==
const double eps = 1e-6;
const double pi = acos(-1);

// 向量 & 点
struct Vect {
    double x, y;

    Vect() = default;
    Vect(double _x, double _y): x(_x), y(_y) {}
    Vect operator+(const Vect &a) const {return Vect(x + a.x, y + a.y);}
    Vect operator-(const Vect &a) const {return Vect(x - a.x, y - a.y);}
    Vect operator*(const double &t) const {return Vect(x * t, y * t);}
    Vect operator/(const double &t) const {return Vect(x / t, y / t);}
    double operator*(const Vect &a) const {return x * a.x + y * a.y;}
    double operator^(const Vect &a) const {return x * a.y - y * a.x;}
    Vect &operator+=(const Vect &a) {return *this = *this + a;}
    Vect &operator-=(const Vect &a) {return *this = *this - a;}
    Vect &operator*=(const double &t) {return *this = *this * t;}
    Vect &operator/=(const double &t) {return *this = *this / t;}
    double length() const {return sqrt(x * x + y * y);}
    Vect normal() const {return *this / length();}
};
using Point = Vect;
// 多边形
struct Polygon: vector<Point> {
    Polygon(initializer_list<Point> &&arg): vector<Point>(arg) {}
    template<typename... argT>
    explicit Polygon(argT &&...args): vector<Point>(forward<argT>(args)...) {}

    // 多边形面积
    double area() const {
        double ans = 0;
        for (int i = 0; i < (int)size(); i++) ans += (*this)[i] ^ (*this)[(i + 1) % size()];
        return ans / 2;
    }
};
// 直线
struct Line {
    // 表示方式为起点 + 向量
    Point s;
    Vect dir;

    Line() = default;
    Line(Point _s, Vect _dir): s(_s), dir(_dir) {}
};
// 线段
struct Segment: Line {
    Segment() = default;
    Segment(Point _s, Vect _dir): Line(_s, _dir) {}
};
// 圆
struct Circle {
    // 表示方式为圆心 + 半径
    Point O;
    double r;

    Circle() = default;
    Circle(Point _O, double _r): O(_O), r(_r) {}
};

// 判断 x < 0 或 x = 0 或 x > 0
int sign(double x) {return fabs(x) <= eps ? 0 : (x > 0 ? 1 : -1);}
// 向量 u 逆时针旋转 rad 后得到的向量
Vect rotate(Vect u, double rad) {
    return Vect(u.x * cos(rad) - u.y * sin(rad), u.x * sin(rad) + u.y * cos(rad));
}
// 两点距离
double dist(Point u, Point v) {return sqrt((u.x - v.x) * (u.x - v.x) + (u.y - v.y) * (u.y - v.y));}
// 两个向量的夹角
double angle(Vect a, Vect b) {return atan2(a ^ b, a * b);}
// 向量 a 到向量 b 的投影
Vect projection(Vect a, Vect b) {return b * (a * b / b.length() / b.length());}
// 点 a 到直线 b 的投影
Point projection(Point a, Line b) {return b.s + projection(a - b.s, b.dir);}
// 点 a 关于直线 b 的轴对称点
Point reflection(Point a, Line b) {return b.s + projection(a - b.s, b.dir) * 2 - a;}
// 判断两条直线是否平行
int isParallel(Line a, Line b) {return !sign(a.dir ^ b.dir);}
// 判断两条直线是否相交
int isIntersecting(Line a, Line b) {return !isParallel(a, b);}
// 判断两条直线是否垂直
int isOrthogonal(Line a, Line b) {return !sign(a.dir * b.dir);}
// 判断点 a 是否在直线 b 上
int isOnLine(Point a, Line b) {return !sign(b.dir ^ (a - b.s));}
// 判断点 a 是否在线段 b 上
int isOnSegment(Point a, Segment b) {
    return isOnLine(a, b) && sign((a - b.s) * (a - b.s - b.dir)) <= 0;
}
// 判断直线 a 与线段 b 是否相交
int isIntersecting(Line a, Segment b) {
    int va = sign((b.s - a.s) ^ a.dir), vb = sign((b.s + b.dir - a.s) ^ a.dir);
    return !va || !vb || va != vb;
}
// 判断两条线段是否相交
int isIntersecting(Segment a, Segment b) {
    return isOnSegment(b.s, a) || isOnSegment(b.s + b.dir, a) ||
        isOnSegment(a.s, b) || isOnSegment(a.s + a.dir, b) ||
        (sign((b.s - a.s) ^ a.dir) != sign((b.s + b.dir - a.s) ^ a.dir) &&
            sign((a.s - b.s) ^ b.dir) != sign((a.s + a.dir - b.s) ^ b.dir));
}
// 直线交点
Point intersection(Line a, Line b) {return b.s + b.dir * (a.dir ^ (b.s - a.s)) / (b.dir ^ a.dir);}
// 点 a 到直线 b 的距离
double dist(Point a, Line b) {return dist(a, projection(a, b));}
// 点 a 到线段 b 的距离
double dist(Point a, Segment b) {
    double ans = min(dist(a, b.s), dist(a, b.s + b.dir));
    Point tmp = projection(a, b);
    if (isOnSegment(tmp, b)) ans = min(ans, dist(a, tmp));
    return ans;
}
// 线段间的距离
double dist(Segment a, Segment b) {
    if (isIntersecting(a, b)) return 0;
    return min({dist(a.s, b), dist(a.s + a.dir, b), dist(b.s, a), dist(b.s + b.dir, a)});
}
// 判断多边形是否是凸包
int isConvex(const Polygon &a) {
    int n = a.size();
    for (int i = 0; i < n; i++)
        if (sign((a[(i + 1) % n] - a[i]) ^ (a[(i + 2) % n] - a[i])) < 0) return 0;
    return 1;
}
// 判断点 a 与多边形 b 的位置关系，返回值：0 - 在多边形外, 1 - 在多边形边上, 2 - 在多边形内部
int Containment(Point a, const Polygon &b) {
    int n = b.size(), val = 0;
    for (int i = 0; i < n; i++) {
        Segment s = Segment(b[i], b[(i + 1) % n] - b[i]);
        if (isOnSegment(a, s)) return 1;
        int t = sign((b[(i + 1) % n] - b[i]) ^ (a - b[i]));
        int va = sign(b[i].y - a.y), vb = sign(b[(i + 1) % n].y - a.y);
        if (t > 0 && va <= 0 && vb > 0) val++;
        if (t < 0 && vb <= 0 && va > 0) val--;
    }
    return val ? 2 : 0;
}
// 点集的凸包
Polygon Andrew(Polygon a) {
    Polygon ans;
    sort(a.begin(), a.end(), [](const Point &u, const Point &v) {
        return u.x != v.x ? u.x < v.x : u.y < v.y;
    });
    ans.push_back(a[0]);
    for (int i = 1; i < (int)a.size(); i++) {
        while (ans.size() > 1) {
            Point tmp = ans.back();
            ans.pop_back();
            if (sign((tmp - ans.back()) ^ (a[i] - tmp)) > 0) {ans.push_back(tmp); break;}
        }
        ans.push_back(a[i]);
    }
    int top = ans.size();
    for (int i = (int)a.size() - 2; i >= 0; i--) {
        while ((int)ans.size() > top) {
            Point tmp = ans.back();
            ans.pop_back();
            if (sign((tmp - ans.back()) ^ (a[i] - tmp)) > 0) {ans.push_back(tmp); break;}
        }
        ans.push_back(a[i]);
    }
    ans.pop_back();
    return ans;
}
// 两个凸包的闵可夫斯基和
Polygon Minkowski(const Polygon &l, const Polygon &r) {
    Polygon diffl, diffr;
    int sizl = l.size(), sizr = r.size();
    for (int i = 1; i < sizl; i++) diffl.push_back(l[i] - l[i - 1]);
    for (int i = 1; i < sizr; i++) diffr.push_back(r[i] - r[i - 1]);
    Polygon ans;
    ans.push_back(l[0] + r[0]);
    int pl = 0, pr = 0;
    while (pl < sizl - 1 && pr < sizr - 1)
        if (sign(diffl[pl] ^ diffr[pr]) < 0) ans.push_back(ans.back() + diffl[pl++]);
        else ans.push_back(ans.back() + diffr[pr++]);
    while (pl < sizl - 1) ans.push_back(ans.back() + diffl[pl++]);
    while (pr < sizr - 1) ans.push_back(ans.back() + diffr[pr++]);
    return ans;
}
// 旋转卡壳求凸包直径
double Diameter(Polygon a) {
    const int n = a.size();
    if (n == 1) return 0;
    if (n == 2) return dist(a[0], a[1]);
    double ans = 0;
    a.push_back(a[0]);
    for (int i = 0, pos = 2; i < n; i++) {
        Point u = a[i], v = a[i + 1];
        ans = max(ans, dist(u, v));
        while (sign(((u - a[pos + 1]) ^ (v - a[pos + 1])) - ((u - a[pos]) ^ (v - a[pos]))) >= 0)
            pos = (pos + 1) % n;
        ans = max({ans, dist(u, a[pos]), dist(v, a[pos])});
    }
    return ans;
}
// 扇形面积
double sector(double r, double rad) {return r * r * rad / 2;}
// 弧度为 rad 的弓形的面积
double arch(double r, double rad) {return sector(r, rad) - r * r * sin(rad) / 2;}
// 两个圆的公切线个数（判断两个圆的位置关系）
int countOfCommonTangentLines(Circle a, Circle b) {
    if (sign(dist(a.O, b.O) - a.r - b.r) > 0) return 4;
    else if (sign(dist(a.O, b.O) - a.r - b.r) == 0) return 3;
    else if (sign(dist(a.O, b.O) - fabs(a.r - b.r)) == 0) return 1;
    else if (sign(dist(a.O, b.O) - fabs(a.r - b.r)) < 0) return 0;
    else return 2;
}
// 两个圆的交点
pair<Point, Point>intersection(Circle a, Circle b) {
    Vect u = b.O - a.O;
    double rad = acos((a.r * a.r + u.length() * u.length() - b.r * b.r) / a.r / u.length() / 2);
    return {rotate(u, -rad) + a.O, rotate(u, rad) + a.O};
}
\end{minted}
%
%% Chapter ? 组合数学
\chapter{组合数学}

\section{组合数的一些有用的公式}

范德蒙德卷积及其推论：

$$
\begin{aligned}
\sum_{i=0}^k \binom{n}{i}\binom{m}{k-i}=\binom{n+m}{k} \\
\sum_{i=-r}^s \binom{n}{r+i}\binom{m}{s-i}=\binom{n+m}{r+s} \\
\sum_{i=1}^n \binom{n}{i}\binom{n}{i-1}=\binom{2n}{n-1} \\
\sum_{i=0}^n \binom{n}{i}^2=\binom{2n}{n} \\
\sum_{i=0}^m \binom{n}{i}\binom{m}{i}=\binom{n+m}{m}
\end{aligned}
$$

用来容斥的东西：

$$\sum_{i=0}^n (-1)^i \binom{n}{i}=[n=0]$$

组合数带权和（对于 $i^k$ 可以得出是关于 $n$ 的 $k$ 次多项式与 $2^{n-k}$ 的乘积）：

$$\sum_{i=0}^n i \binom{n}{i}=n2^{n-1}$$

$$\sum_{i=0}^n i^2 \binom{n}{i}=n(n+1)2^{n-2}$$

上指标求和：

$$\sum_{i=m}^n \binom{i}{m}=\binom{n+1}{m+1}$$

三项式版恒等式：

$$\binom{n}{m}\binom{m}{k}=\binom{n}{k}\binom{n-k}{m-k}$$

平行恒等式：

$$\sum_{i=0}^n \binom{m+i}{i}=\binom{m+n+1}{n}$$

上指标卷积：

$$\sum_{i=0}^n \binom{i}{a}\binom{n-i}{b}=\binom{n+1}{a+b+1}$$

二项式反演：

$$f_n=\sum_{i=0}^n \binom{n}{i}g_i \leftrightarrow g_n=\sum_{i=0}^n (-1)^{n-i} \binom{n}{i} f_i$$

$$f_n=\sum_{i=n}^m \binom{i}{n} g_i \leftrightarrow g_n=\sum_{i=n}^m (-1)^{i-n}\binom{i}{n}f_i$$

多项式定理：

设 $n$ 为正整数，$x_i$ 为实数，则：

$$(x_1+x_2+...+x_t)^n=\sum_{\text{满足 $n_1+n_2+...+n_t=n$ 的非负整数解}} \binom{n}{n_1,n_2,...,n_t}x_1^{n_1}x_2^{n_2}\dots x_t^{n_t}$$

其中的 $\binom{n}{n_1,n_2,...,n_t}$ 是多重集的排列数，取 $x_1=x_2=...=x_t=1$，得到 $\sum \binom{n}{n_1,n_2,...,n_t}=t^n$。

错位排列：

$$D_n=\sum_{k=0}^n(-1)^k\binom{n}{k}(n-k)!$$

卡特兰数：

$$H_n=\frac{1}{n+1}\binom{2n}{n}=\frac{(2n)!}{n!(n+1)!}$$

$$H_n=\binom{2n}{n}-\binom{2n}{n-1}$$

$$H_n=\frac{H_{n-1}(4n-2)}{n+1}$$



\section{分拆数}

$1,1,2,3,5,7,11,15,22,30,42,56,77,101,135,176,231,297,385,490,627,792,1002,1255,1575,1958,$

$2436,3010,3718,4565,5604,6842,8349,10143,12310,14883,17977,21637,26015,31185,37338,44583,$

$53174,63261,75175,89134,105558,124754,147273,173525,204226,239943,281589,329931,386155,$

$451276,526823,614154,715220,831820,966467,1121505,1300156,1505499,1741630,2012558,2323520,$

$2679689,3087735,3554345,4087968,4697205,5392783,6185689,7089500,8118264,9289091,10619863,$

$12132164,13848650,15796476,18004327,20506255,23338469,26543660,30167357,34262962,38887673,$

$44108109,49995925,56634173,64112359,72533807,82010177,92669720,104651419,$

$118114304,133230930,150198136,169229875,190569292$

%
%% Chapter ? 数学杂项
\chapter{数学杂项}

\section{群论}

\section{线性规划}

\begin{tcolorbox}
\textbf{线性规划标准型}：

$$ \max \sum_{i=1}^n c_ix_i $$
$$\text{s.t.\ \ \ } \sum_{j=1}^n a_{i,j}x_j \le b_i , i=1,2,3,\dots,m $$
$$x_i\ge 0,i=1,2,3,\dots,n$$


要求最大化一个关于这些变量的线性函数，满足每个限制的形式是一个线性函数小于等于某个值，且要求所有变量非负。
\end{tcolorbox}

写成矩阵的形式：

$$ \max \mathbf{C}^T \mathbf{x} $$
$$ \text{s.t.\ \ \ } \mathbf{Ax} \le \mathbf{B} $$
$$ \mathbf{x \ge 0}$$

\subsection{线性规划对偶}

对于上述标准型线性规划，其对偶形式是：

$$ \min \mathbf{B}^T \mathbf{y} $$
$$ \text{s.t.\ \ \ } \mathbf{A^Ty} \le \mathbf{C} $$
$$ \mathbf{y \ge 0}$$

可以证明，两个互为对偶的线性规划的最优解相等。

\subsection{单纯形法}

以下代码可以求解一个标准型线性规划，数据范围 $1\leq n,m \leq 20$，$0 \leq |a_{ij}|,|b_i|,|c_j|\leq 100$。

\begin{minted}{c++}
typedef long double ld;
typedef __float128 f128;
const int N = 25;
const f128 EPS = 1e-16L , INF = 1e50L;
mt19937 rnd(114514);
f128 fabs(const f128 &x){return x < 0 ? -x : x;}
namespace LP
{
	const int N = ::N , M = ::N;
	#define b(i) (a[i][0])
	#define c(i) (a[0][i])
	int n , m , id[N + M];
	f128 a[M][N] , ans[N + M];
	void Pivot(int l , int e)
	{
		f128 k = a[l][e]; a[l][e] = 1;
		swap(id[e] , id[l + n]);
		for(int i = 0 ; i <= n ; i++)
			a[l][i] /= k;
		for(int i = 0 ; i <= m ; i++)
		{
			if(i == l || fabs(a[i][e]) < EPS)continue ;
			f128 k = a[i][e]; a[i][e] = 0;
			b(i) += (i ? -1 : 1) * k * b(l);
			for(int j = 1 ; j <= n ; j++)
				a[i][j] -= k * a[l][j];
		}
	}
	bool Simplex()
	{
		while(1)
		{
			int l = 0 , e = 0; f128 mn = INF;
			for(int i = 1 ; i <= n && !e ; i++)
				if(a[0][i] > EPS)e = i;
			if(!e)break ;
			for(int i = 1 ; i <= m ; i++)
			{
				f128 v = b(i) / a[i][e];
				if(a[i][e] > EPS && (v + EPS < mn || 
				(fabs(mn - v) < EPS && id[i + n] < id[l + n])))
					mn = v , l = i;
			}
			if(!l)return 1;
			Pivot(l , e);
		}
		return 0;
	}
	bool Init(int nn , int mm)
	{
		n = nn , m = mm , b(0) = 0;
		iota(id + 1 , id + n + m + 1 , 1);
		while(1)
		{
			int l = 0 , e = 0;
			for(int i = 1 ; i <= m ; i++)
				if(b(i) < -EPS && (!l || rnd() & 1))l = i;
			if(!l)break ;
			for(int i = 1 ; i <= n ; i++)
				if(a[l][i] < -EPS && (!e || rnd() & 1))e = i;
			if(!e)return 1;
			Pivot(l , e);
		}
		return 0;
	}
}
using LP::a;
using LP::ans;
using LP::id;
signed main()
{
	ios::sync_with_stdio(0);
	cin.tie(0) , cout.tie(0);
	int n , m , tp; ld x; bool ex48;
	cin >> n >> m >> tp;
	for(int i = 1 ; i <= n ; i++)//c
		cin >> x , a[0][i] = x;
	for(int i = 1 ; i <= m ; i++)
	{
		for(int j = 1 ; j <= n ; j++)//a
			cin >> x , a[i][j] = x;
		cin >> x , a[i][0] = x;//b
	}
	if(LP::Init(n , m))cout << "Infeasible\n";
	else if(LP::Simplex())cout << "Unbounded\n";
	else
	{
		cout << setprecision(8) << fixed << (ld)a[0][0] << "\n";
		if(tp == 1)
		{
			memset(ans , 0 , sizeof ans);
			for(int i = 1 ; i <= m ; i++)ans[id[i + n]] = b(i);
			for(int i = 1 ; i <= n ; i++)cout << (ld)ans[i] << " \n"[i == n];
		}
	}
	return 0;
}
\end{minted}

\section{微积分}

\section{高精度}

\begin{minted}{c++}
// == Main ==
// if use read or print, don't close ios_base::sync_with_stdio
struct ubigint : public vector<int> {
    const static int BASE = (int)1e8;
    const static int WS = 8;  //与BASE一起改
    ubigint() {}
    ubigint(int t) : vector<int>(1, t) {}
    ubigint(string s) {
        for (int i = s.size() - 1, j, t; i >= 0; i -= WS) {
            for (j = max(0, i - WS + 1), t = 0; j <= i; ++j) t = (t * 10 + (s[j] - '0'));
            emplace_back(t);
        }
    }
    unsigned int digits() const {
        if (isZero()) return 1;
        unsigned int d = (size() - 1u) * WS;
        int x = back();
        while (x) ++d, x /= 10;
        return d;
    }
    const bool isZero() const { return empty() or (size() == 1u and !front()); }
    int const& operator[](const int n) const { return (n < size()) ? *(cbegin() + n) : (int const&)0; }
    int& operator[](const int n) { return *(begin() + n); }
    operator string() {
        if (isZero()) return "0";
        string ans;
        stringstream ss;
        ss << back();
        for (int i = (int)size() - 2; ~i; --i) ss << setw(WS) << setfill('0') << *(begin() + i);
        ss >> ans;
        return ans;
    }
    friend void read(ubigint& a) {
        string s;
        char c = getchar();
        while (!isdigit(c)) c = getchar();
        while (isdigit(c)) s.push_back(c), c = getchar();
        a = ubigint(s);
    }
    friend void print(const ubigint& a) {
        if (a.isZero()) return (void)putchar('0');
        printf("%d", a.back());
        for (int i = (int)a.size() - 2; ~i; --i) printf("%0*d", WS, a[i]);
    }
    friend bool operator==(ubigint const& a, ubigint const& b) {
        if (a.size() != b.size()) return 0;
        for (int i = 0; i < (int)a.size(); ++i)
            if (a[i] != b[i]) return 0;
        return 1;
    }
    friend bool operator!=(ubigint const& a, ubigint const& b) { return !(a == b); }
    friend bool operator<(ubigint const& a, ubigint const& b) {
        if (a.size() != b.size()) return a.size() < b.size();
        for (int i = (int)a.size() - 1; ~i; --i)
            if (a[i] != b[i]) return a[i] < b[i];
        return 0;
    }
    friend bool operator>(ubigint const& a, ubigint const& b) { return b < a; }
    friend bool operator<=(ubigint const& a, ubigint const& b) {
        if (a.size() != b.size()) return a.size() < b.size();
        for (int i = (int)a.size() - 1; ~i; --i)
            if (a[i] != b[i]) return a[i] < b[i];
        return 1;
    }
    friend bool operator>=(ubigint const& a, ubigint const& b) { return b <= a; }
    friend ubigint operator+(ubigint const& a, ubigint const& b) {
        ubigint c = a;
        c.resize(max(a.size(), b.size()) + 1);
        for (int i = 0, en = b.size(); i < en; ++i) {
            c[i] += b[i];
            if (c[i] >= BASE) c[i] -= BASE, ++c[i + 1];
        }
        for (int i = b.size(), en = c.size() - 1; i < en; ++i)
            if (c[i] >= BASE) c[i] -= BASE, ++c[i + 1];
        if (!c.back()) c.pop_back();
        return c;
    }
    friend ubigint operator-(ubigint const& a, ubigint const& b) {
        ubigint c = a;
        for (int i = 0, en = b.size(); i < en; ++i) {
            c[i] -= b[i];
            if (c[i] < 0) c[i] += BASE, --c[i + 1];
        }
        for (int i = b.size(), en = c.size(); i < en; ++i) {
            if (c[i] < 0)
                c[i] += BASE, --c[i + 1];
            else
                break;
        }
        while (c.size() > 1u && !c.back()) c.pop_back();
        return c;
    }
    friend ubigint operator*(ubigint const& a, int o) {
        if (!o or a.isZero()) return 0;
        int n = a.size();
        ubigint c;
        c.resize(n);
        long long carry = 0;
        for (int i = 0; i < n; ++i) {
            carry += (long long)a[i] * o;
            c[i] = carry % BASE;
            carry /= BASE;
        }
        if (carry) c.emplace_back(carry);
        return c;
    }
    friend ubigint operator*(ubigint const& a, ubigint const& b) {
        if (a.isZero() || b.isZero()) return 0;
        int sa = a.size(), sb = b.size();
        vector<long long> t(sa + sb);
        for (int i = 0; i < sa; ++i)
            for (int j = 0; j < sb; ++j) t[i + j] += (long long)a[i] * b[j];
        for (int i = 0, en = t.size() - 1; i < en; ++i) t[i + 1] += t[i] / BASE, t[i] %= BASE;
        if (!t.back()) t.pop_back();
        ubigint c;
        c.resize(t.size());
        for (int i = 0, en = t.size(); i < en; ++i) c[i] = (int)t[i];
        return c;
    }
    friend ubigint operator/(ubigint const& a, ubigint const& b) {
        assert(!b.isZero());
        if (a.size() < b.size()) return 0;
        ubigint c, d;
        c.assign(a.end() - b.size() + 1, a.end());
        for (int i = a.size() - b.size(); ~i; --i) {
            c.insert(c.begin(), *(a.begin() + i));
            long long t = (c.size() > b.size()) ? ((long long)c.back() * BASE + *(c.crbegin() + 1)) : c.back();
            int l = (t / (b.back() + 1)), r = ((t + 1) / b.back()), mid;
            while (l < r) {
                mid = (l + r + 1) >> 1;
                if (b * mid <= c)
                    l = mid;
                else
                    r = mid - 1;
            }
            c -= b * l;
            if (!c.back()) c.pop_back();
            d.emplace_back(l);
        }
        reverse(d.begin(), d.end());
        if (d.size() > 1u && !d.back()) d.pop_back();
        return d;
    }
    friend ubigint divide(ubigint const& a, ubigint const& b, ubigint& c) {
        assert(!b.isZero());
        if (a.size() < b.size() or a.isZero()) return c = a, 0;
        ubigint d;
        c.assign(a.end() - b.size() + 1, a.end());
        for (int i = a.size() - b.size(); ~i; --i) {
            c.insert(c.begin(), *(a.begin() + i));
            long long t = (c.size() > b.size()) ? ((long long)c.back() * BASE + *(c.crbegin() + 1)) : c.back();
            int l = (t / (b.back() + 1)), r = ((t + 1) / b.back()), mid;
            while (l < r) {
                mid = (l + r + 1) >> 1;
                if (b * mid <= c)
                    l = mid;
                else
                    r = mid - 1;
            }
            c -= b * l;
            if (!c.back()) c.pop_back();
            d.emplace_back(l);
        }
        reverse(d.begin(), d.end());
        if (d.size() > 1u && !d.back()) d.pop_back();
        return d;
    }
    friend ubigint operator ^ (ubigint x , int p)
    {
        ubigint res = 1;
        for(; p > 0 ; p >>= 1 , x = x * x)
            if(p & 1)res = res * x;
        return res;
    }   
    inline friend ubigint sqrt(const ubigint& n , const int& m)
    {
        ubigint l = 0 , r = 1;
        if(m == 1)return n;
        while((r ^ m) <= n)r = r * 2;
        l = r / 2;
        while(l < r)
        {
            ubigint mid = (l + r + 1) / 2;
            if((mid ^ m) <= n)l = mid;
            else r = mid - 1;
        }
        return l;
    }
    friend ubigint operator%(ubigint const& a, ubigint const& b) { return a - a / b * b; }
    friend ubigint const& operator+=(ubigint& a, ubigint const& b) { return a = a + b; }
    friend ubigint const& operator-=(ubigint& a, ubigint const& b) { return a = a - b; }
    friend ubigint const& operator*=(ubigint& a, int const& b) { return a = a * b; }
    friend ubigint const& operator*=(ubigint& a, ubigint const& b) { return a = a * b; }
    friend ubigint const& operator/=(ubigint& a, ubigint const& b) { return a = a / b; }
    friend ubigint const& operator%=(ubigint& a, ubigint const& b) { return a = a % b; }
};
struct bigint {
    bool f;
    ubigint t;
    bigint() : f(0) {}
    bigint(int t) : f(t < 0), t(ubigint(abs(t))) {}
    bigint(bool f, ubigint t) : f(f), t(t) {}
    unsigned int digits() const { return t.digits(); }
    const bool isZero() const { return t.isZero(); }
    int const& operator[](const int n) const { return (n < t.size()) ? *(t.cbegin() + n) : (int const&)0; }
    int& operator[](const int n) { return *(t.begin() + n); }
    operator string() {
        if (t.isZero()) return "0";
        string ans;
        stringstream ss;
        if (f) ss << '-';
        ss << t.back();
        for (int i = (int)t.size() - 2; ~i; --i) ss << setw(t.WS) << setfill('0') << t[i];
        ss >> ans;
        return ans;
    }
    friend void read(bigint& a) {
        a.f = 0;
        string s;
        char c = getchar();
        for (; !isdigit(c); c = getchar())
            if (c == '-') a.f = 1;
        while (isdigit(c)) s.push_back(c), c = getchar();
        a.t = ubigint(s);
    }
    friend void print(const bigint& a) {
        if (a.t.isZero()) return (void)putchar('0');
        if (a.f) putchar('-');
        print(a.t);
    }
    friend bigint abs(bigint const& a) { return bigint(0, a.t); }
    friend bool operator==(bigint const& a, bigint const& b) { return (a.f == b.f) && (a.t == b.t); }
    friend bool operator!=(bigint const& a, bigint const& b) { return !(a == b); }
    friend bool operator<(bigint const& a, bigint const& b) {
        if (a.f != b.f) return a.f;
        return a.f ? a.t > b.t : a.t < b.t;
    }
    friend bool operator>(bigint const& a, bigint const& b) { return b < a; }
    friend bool operator<=(bigint const& a, bigint const& b) {
        if (a.f != b.f) return a.f;
        return a.f ? a.t >= b.t : a.t <= b.t;
    }
    friend bool operator>=(bigint const& a, bigint const& b) { return b <= a; }
    friend bigint operator-(bigint const& a) { return bigint(!a.f, a.t); }
    friend bigint operator+(bigint const& a, bigint const& b) {
        if (a.f == b.f) return bigint(a.f, a.t + b.t);
        if (a.t > b.t) return bigint(a.f, a.t - b.t);
        if (a.t < b.t) return bigint(b.f, b.t - a.t);
        return 0;
    }
    friend bigint operator-(bigint const& a, bigint const& b) {
        if (a.f == b.f) {
            if (a.t > b.t) return bigint(a.f, a.t - b.t);
            if (a.t < b.t) return bigint(!a.f, b.t - a.t);
            return 0;
        }
        return bigint(a.f, a.t + b.t);
    }
    friend bigint operator*(bigint const& a, bigint const& b) { return bigint(a.f ^ b.f, a.t * b.t); }
    friend bigint operator*(bigint const& a, int const& b) { return bigint(a.f ^ (b < 0), a.t * b); }
    friend bigint divide(bigint const& a, bigint const& b, bigint& c) { return bigint(a.f ^ b.f, divide(a.t, b.t, (c.f = a.f, c.t))); }
    friend bigint operator/(bigint const& a, bigint const& b) { return bigint(a.f ^ b.f, a.t / b.t); }
    friend bigint operator%(bigint const& a, bigint const& b) { return bigint(a.f, a.t % b.t); }
    friend bigint const& operator+=(bigint& a, bigint const& b) { return a = a + b; }
    friend bigint const& operator-=(bigint& a, bigint const& b) { return a = a - b; }
    friend bigint const& operator*=(bigint& a, int const& b) { return a = a * b; }
    friend bigint const& operator*=(bigint& a, bigint const& b) { return a = a * b; }
    friend bigint const& operator/=(bigint& a, bigint const& b) { return a = a / b; }
    friend bigint const& operator%=(bigint& a, bigint const& b) { return a = a % b; }
};
\end{minted}

%
%% Chapter ? 杂项
\chapter{杂项}

\section{取模类}

\begin{minted}{c++}
// == Main ==
struct mint {
    static const int mod = 998244353;
    int v;

    mint() = default;
    mint(int _v): v((_v % mod + mod) % mod) {}
    explicit operator int() const {return v;}
    mint operator+(const mint &x) const {return v + x.v - (v + x.v < mod ? 0 : mod);}
    mint &operator+=(const mint &x) {return *this = *this + x;}
    mint operator-(const mint &x) const {return v - x.v + (v - x.v >= 0 ? 0 : mod);}
    mint &operator-=(const mint &x) {return *this = *this - x;}
    mint operator*(const mint &x) const {return (long long)v * x.v % mod;}
    mint &operator*=(const mint &x) {return *this = *this * x;}
    mint inv() const {
        mint a(*this), ans(1);
        int b(mod - 2);
        while (b) {
            if (b & 1) ans *= a;
            a *= a;
            b >>= 1;
        }
        return ans;
    }
    mint operator/(const mint &x) const {return *this * x.inv();}
    mint &operator/=(const mint &x) {return *this = *this / x;}
    mint operator-() {return mint(-v);}
};
\end{minted}

\subsection{Barrett 约减}

当模数不固定时可以加速。

用法：在构造函数中传模数，使用方法为 \lstinline|F.reduce(x)|，其中 $x$ 是需要取模的数。

\begin{minted}{c++}
// == Main ==
struct Barrett {
    unsigned long long b, m;
    Barrett(unsigned long long b = 2): b(b), m((__uint128_t(1) << 64) / b) {}
    unsigned long long reduce(long long x) {
        unsigned long long r = (__uint128_t(x + b) * m) >> 64;
        unsigned long long q = (x + b) - b * r;
        return q >= b ? q - b : q;
    }
} F;
\end{minted}

\section{对拍脚本}

\begin{minted}{bash}
#!/usr/bin/bash

declare -i num=0

while [ true ]; do
    ./mkdata > in.txt
    time ./mine < in.txt > out.txt
    ./correct < in.txt > ans.txt
    diff out.txt ans.txt
    if [ $? -ne 0 ]; then
        echo "WA"
        break
    fi
    num=num+1
    echo "Passed $num tests."
done
\end{minted}

\section{VS Code 配置}

\subsection{User Tasks}

\begin{minted}{json}
{
    // See https://go.microsoft.com/fwlink/?LinkId=733558
    // for the documentation about the tasks.json format
    "version": "2.0.0",
    "tasks": [
        {
            "type": "shell",
            "label": "My C++ Runner",
            "detail": "Build and Run Current C++ Program",
            "command": [ // 三个编译方式保留一个即可。
                "clear",
                "&&",
                "g++ ${file} -o ${fileDirname}/${fileBasenameNoExtension} -std=c++14 -Wall -Wextra && echo '== Normal =='",
                "g++ ${file} -o ${fileDirname}/${fileBasenameNoExtension} -std=c++14 -Wall -Wextra -O2 && echo '== O2 =='",
                "g++ ${file} -o ${fileDirname}/${fileBasenameNoExtension} -std=c++14 -Wall -Wextra -fsanitize=undefined,address && echo '== UB Check =='",
                "&&",
                "gnome-terminal -- bash -c \"ulimit -s 524288; time ${fileDirname}/${fileBasenameNoExtension}; read -p 'Press ENTER to continue...'; exit\""
            ],
            "problemMatcher": [ // 非必要
                "$gcc"
            ],
            "group": { // 非必要
                "kind": "build",
                "isDefault": true
            },
            "presentation": { // 非必要
                "showReuseMessage": false
            }
        }
    ]
}
\end{minted}

\subsection{设置}

\begin{itemize}
    \item 字体大小：16。（\lstinline|"editor.fontSize": 16|）
    \item 添加多个光标的方式：ctrl。（\lstinline|"editor.multiCursorModifier": "ctrlCmd"|）
    \item 不适用空格代替 Tab。（\lstinline|"editor.insertSpaces": false|）
    \item 不允许 Enter 进行代码补全。（\lstinline|"editor.acceptSuggestionOnEnter": "off"|）
    \item 标尺：110。（\lstinline|"editor.rulers": [110]|）
    \item 平滑。（\lstinline|"editor.cursorSmoothCaretAnimation": "on"|）
    \item 标题栏外观。（\lstinline|"window.titleBarStyle": "custom"|）
\end{itemize}

totally：

\begin{minted}{json}
{
    "editor.fontSize": 16,
    "editor.multiCursorModifier": "ctrlCmd",
    "editor.insertSpaces": false,
    "editor.acceptSuggestionOnEnter": "off",
    "editor.rulers": [110],
    "editor.cursorSmoothCaretAnimation": "on",
    "window.commandCenter": false
}
\end{minted}

\subsection{快捷键}

\begin{itemize}
    \item 切换块注释：Ctrl+Shift+A -> Ctrl+Shift+/
    \item 运行任务：Ctrl+Shift+B -> F11
    \item 向上移动行：Alt+up -> Ctrl+Shift+up
    \item 向下移动行：Alt+down -> Ctrl+Shift+down
\end{itemize}

\end{document}
